\section{School of Computing}\label{school-of-computing}

\subsection{FACULTY OF ENGINEERING AND PHYSICAL
SCIENCES}\label{faculty-of-engineering-and-physical-sciences}

\section{SKYWAVE: A Machine Learning Framework for Predicting High
Frequency Radio Reception
Probability}\label{skywave-a-machine-learning-framework-for-predicting-high-frequency-radio-reception-probability}

\textbf{{[}Kuppa Ganesh{]}}

Submitted in accordance with the requirements for the degree of\\
\textbf{MSc Advanced Computer Science (Data Analytics)}\\
2025/2026

The candidate confirms that the following have been submitted:

{\def\LTcaptype{none} % do not increment counter
\begin{longtable}[]{@{}
  >{\raggedright\arraybackslash}p{(\linewidth - 4\tabcolsep) * \real{0.3333}}
  >{\raggedright\arraybackslash}p{(\linewidth - 4\tabcolsep) * \real{0.3333}}
  >{\raggedright\arraybackslash}p{(\linewidth - 4\tabcolsep) * \real{0.3333}}@{}}
\toprule\noalign{}
\begin{minipage}[b]{\linewidth}\raggedright
Items
\end{minipage} & \begin{minipage}[b]{\linewidth}\raggedright
Format
\end{minipage} & \begin{minipage}[b]{\linewidth}\raggedright
Recipient(s) and Date
\end{minipage} \\
\midrule\noalign{}
\endhead
\bottomrule\noalign{}
\endlastfoot
Deliverable 1 & Project Report & SSO (Date of Submission) \\
Deliverable 2 & Source Code Repository (URL) & Supervisor, Assessor
(Date of Submission) \\
\end{longtable}
}

\textbf{Type of Project:} Empirical Investigation

The candidate confirms that the work submitted is their own and the
appropriate credit has been given where reference has been made to the
work of others.\\
I understand that failure to attribute material which is obtained from
another source may be considered as plagiarism.

(Signature of
student)\_\_\_\_\_\_\_\_\_\_\_\_\_\_\_\_\_\_\_\_\_\_\_\_\_\_

© 2026 The University of Leeds and Kuppa Ganesh

\begin{center}\rule{0.5\linewidth}{0.5pt}\end{center}

\subsection{Summary}\label{summary}

High Frequency (HF) radio communication remains a critical
infrastructure for long-distance data exchange in environments where
conventional telecommunications are unavailable or unreliable. HF
signals propagate by refracting off ionised layers in the upper
atmosphere, a process governed by complex interactions between solar
radiation, time of day, and seasonal cycles. Because the ionosphere is
inherently dynamic, the success of a transmission cannot be determined
deterministically from transmitter settings alone. The aim of this
project is to develop and evaluate a machine learning framework, termed
SKYWAVE, that estimates the probability of successful reception for HF
transmissions using historical operational logs.

This investigation utilises a dataset comprising approximately 7,700
transmission records collected over a six-month period. The data
presents a class imbalance, with successful receptions occurring in
roughly 48\% of attempts. To address this, the framework employs a
Random Forest classifier optimised with balanced class weights and
isotonic probability calibration. A significant methodological
contribution of this project is the implementation of a rigorous leakage
prevention protocol, which explicitly excludes post-transmission metrics
from the feature space, ensuring that the model learns propagation
patterns rather than memorising dataset artefacts.

Feature engineering incorporates physics-informed interaction terms,
such as the dependency of higher frequency bands on daylight conditions.
The model is evaluated using a dual-validation strategy comprising
stratified k-fold cross-validation and chronological splitting to assess
temporal generalisation. The final model achieves a Recall of
approximately 0.85, a Precision of 0.91, and a Brier Score of 0.097 on
the held-out test set. These results demonstrate that pre-transmission
parameters---specifically antenna tuning, carrier frequency, and
temporal indicators---can reliably predict reception likelihood. The
calibrated probability outputs enable radio operators to make risk-aware
scheduling decisions, conserving power and improving communication
efficiency in remote operational contexts.

\begin{center}\rule{0.5\linewidth}{0.5pt}\end{center}

\subsection{Acknowledgements}\label{acknowledgements}

I would first like to express my sincere gratitude to my supervisor,
{[}Supervisor Name{]}, for their continuous guidance, support, and
encouragement throughout this project. Their insightful recommendations
and feedback were crucial from the very beginning, particularly in
refining the methodology for leakage prevention and feature engineering.

I am equally grateful to my assessor, {[}Assessor Name{]}, for providing
helpful guidance during the assessment meetings, which significantly
contributed to the structure and academic rigour of this project.

Finally, I would like to thank my family and colleagues for their
patience and support during the intensive research and writing phases of
this degree.

\begin{center}\rule{0.5\linewidth}{0.5pt}\end{center}

\section{Table of Contents}\label{table-of-contents}

\begin{itemize}
\tightlist
\item
  Chapter 1. Introduction

  \begin{itemize}
  \tightlist
  \item
    1.1 Project Aim
  \item
    1.2 Objectives
  \item
    1.3 Deliverables
  \item
    1.4 Ethical, Legal, and Social Issues
  \end{itemize}
\item
  Chapter 2. Background Research

  \begin{itemize}
  \tightlist
  \item
    2.1 Literature Survey

    \begin{itemize}
    \tightlist
    \item
      2.1.1 Ionospheric Propagation and HF Reception Dynamics
    \item
      2.1.2 Machine Learning Applications in Radio Communication
    \item
      2.1.3 Imbalanced Classification and Probability Calibration
    \item
      2.1.4 Temporal Validation and Dataset Shift in Non-Stationary
      Domains
    \end{itemize}
  \item
    2.2 Methods and Techniques

    \begin{itemize}
    \tightlist
    \item
      2.2.1 Data Preprocessing and Leakage Prevention
    \item
      2.2.2 Physics-Informed Feature Engineering
    \item
      2.2.3 Model Architectures for Tabular Prediction
    \item
      2.2.4 Probability Calibration and Metric Selection
    \end{itemize}
  \item
    2.3 Choice of Methods

    \begin{itemize}
    \tightlist
    \item
      2.3.1 Algorithmic Selection and Optimisation Strategy
    \item
      2.3.2 Validation Protocol and Evaluation Framework
    \end{itemize}
  \end{itemize}
\item
  Chapter 3. Datasets and Experimental Design

  \begin{itemize}
  \tightlist
  \item
    3.1 Datasets and Data Sources

    \begin{itemize}
    \tightlist
    \item
      3.1.1 Data Collection and Source Description
    \item
      3.1.2 Variable Taxonomy and Operational Mapping
    \item
      3.1.3 Data Quality Assessment and Preprocessing Protocol
    \end{itemize}
  \item
    3.2 Experimental Design and Protocol

    \begin{itemize}
    \tightlist
    \item
      3.2.1 System Architecture and Pipeline Overview
    \item
      3.2.2 Feature Engineering Strategy
    \item
      3.2.3 Model Selection and Training Protocol
    \item
      3.2.4 Validation Framework and Evaluation Metrics
    \end{itemize}
  \end{itemize}
\item
  Chapter 4. Results of the Empirical Investigation

  \begin{itemize}
  \tightlist
  \item
    4.1 Data Preparation and Feature Engineering Outcomes

    \begin{itemize}
    \tightlist
    \item
      4.1.1 Dataset Cleaning and Quality Assurance
    \item
      4.1.2 Temporal Alignment and Chronological Integrity
    \item
      4.1.3 Physics-Informed Feature Composition
    \item
      4.1.4 Multicollinearity and Feature Selection Validation
    \end{itemize}
  \item
    4.2 Model Training and Probability Calibration

    \begin{itemize}
    \tightlist
    \item
      4.2.1 Algorithm Configuration and Hyperparameter Selection
    \item
      4.2.2 Class Imbalance Mitigation Strategy
    \item
      4.2.3 Isotonic Probability Calibration Protocol
    \item
      4.2.4 Training Dynamics and Convergence Analysis
    \end{itemize}
  \item
    4.3 Performance Metrics and Validation Comparison

    \begin{itemize}
    \tightlist
    \item
      4.3.1 Metric Definitions and Operational Interpretation
    \item
      4.3.2 Stratified Cross-Validation Results
    \item
      4.3.3 Chronological Split Validation
    \item
      4.3.4 Generalisation Gap and Temporal Drift Analysis
    \item
      4.3.5 Statistical Significance and Confidence Intervals
    \end{itemize}
  \item
    4.4 Feature Importance and Domain Alignment

    \begin{itemize}
    \tightlist
    \item
      4.4.1 Global Feature Importance Rankings
    \item
      4.4.2 Physical Validation of Dominant Predictors
    \item
      4.4.3 Ablation Studies and Feature Group Contribution
    \item
      4.4.4 Interaction Effects and Non-Linear Dependencies
    \end{itemize}
  \item
    4.5 Robustness Analysis and Error Characterization

    \begin{itemize}
    \tightlist
    \item
      4.5.1 Confusion Matrix and Error Taxonomy
    \item
      4.5.2 Failure Mode Analysis and Operational Context
    \item
      4.5.3 Sensitivity Analysis and Parameter Perturbation
    \end{itemize}
  \item
    4.6 Computational Efficiency and Resource Utilisation

    \begin{itemize}
    \tightlist
    \item
      4.6.1 Inference Latency and Throughput
    \item
      4.6.2 Memory Footprint and Scalability
    \item
      4.6.3 Calibration Overhead
    \end{itemize}
  \item
    4.7 Chapter Summary
  \end{itemize}
\item
  Chapter 5. Validation of Results

  \begin{itemize}
  \tightlist
  \item
    5.1 Technical Evaluation and Robustness Assessment

    \begin{itemize}
    \tightlist
    \item
      5.1.1 Leakage Prevention Verification
    \item
      5.1.2 Calibration Validation and Threshold Robustness
    \item
      5.1.3 Temporal Generalisation Validation
    \end{itemize}
  \item
    5.2 Hypothesis Verification and Objective Alignment
  \item
    5.3 Error Analysis and Operational Interpretation

    \begin{itemize}
    \tightlist
    \item
      5.3.1 Confusion Matrix Breakdown
    \item
      5.3.2 False Negative Analysis
    \item
      5.3.3 False Positive Analysis
    \item
      5.3.4 Threshold Optimisation for Deployment
    \end{itemize}
  \item
    5.4 Validation Against Domain Knowledge
  \item
    5.5 Validation Constraints and Limitations
  \item
    5.6 Validation Outcome Summary
  \end{itemize}
\item
  Chapter 6. Conclusions and Future Work

  \begin{itemize}
  \tightlist
  \item
    6.1 Conclusions

    \begin{itemize}
    \tightlist
    \item
      6.1.1 Summary of Research Findings
    \item
      6.1.2 Achievement of Project Objectives
    \item
      6.1.3 Operational and Academic Significance
    \end{itemize}
  \item
    6.2 Future Work

    \begin{itemize}
    \tightlist
    \item
      6.2.1 Addressing Dataset Limitations and Temporal Drift
    \item
      6.2.2 Integrating Real-Time Atmospheric Indices
    \item
      6.2.3 Advanced Model Architectures and Multi-Task Learning
    \item
      6.2.4 Deployment Infrastructure and Edge Computing Integration
    \item
      6.2.5 Operator Decision-Support Interface and Threshold
      Optimisation
    \item
      6.2.6 Broader Research Trajectories
    \end{itemize}
  \end{itemize}
\item
  List of References
\item
  Appendices

  \begin{itemize}
  \item
    Appendix A: External Materials
  \item
    Appendix B: Ethical Issues Addressed
  \item ~
    \subsection{Appendix C: Supplementary Figures and
    Tables}\label{appendix-c-supplementary-figures-and-tables}
  \end{itemize}
\end{itemize}

\section{Chapter 1. Introduction}\label{chapter-1.-introduction}

This chapter is divided into four sections which provide an overview of
the project's aim in Section 1.1, the objectives in Section 1.2,
deliverables specified in Section 1.3, and the ethical, legal, and
societal challenges addressed throughout the project in Section 1.4.
More specifically, these sections describe the project's goals,
particular tasks to be accomplished, expected outcomes, and concerns
regarding privacy and social impact.

\subsection{1.1 Project Aim}\label{project-aim}

The aim of this project is to develop and evaluate a machine learning
framework that estimates the probability of successful reception for
High Frequency (HF) radio transmissions using historical operational
logs.

HF radio communication remains a critical medium for long-distance data
exchange in environments where conventional telecommunications
infrastructure is unavailable, degraded, or intentionally restricted.
Signals in the HF band (3--30 MHz) propagate by refracting off ionised
layers in the upper atmosphere, a process governed by complex
interactions between solar radiation, geomagnetic activity, time of day,
and seasonal cycles. Because the ionosphere is inherently dynamic, the
success of a transmission cannot be determined deterministically from
transmitter settings alone. Instead, reception outcomes exhibit
stochastic behaviour that depends on the instantaneous state of the
propagation channel. This project addresses the operational challenge of
predicting reception likelihood by analysing pre-transmission
parameters---such as carrier frequency, antenna standing wave ratio
(SWR), transmitted power, and temporal conditions---to identify patterns
that correlate with successful decoding at remote gateway stations.

In operational contexts, radio practitioners currently rely on static
propagation charts, empirical band plans, or heuristic scheduling rules
that do not account for real-time channel variability. These approaches
often result in repeated failed attempts, unnecessary power expenditure,
and inefficient spectrum utilisation. To address this limitation, the
framework developed in this study produces calibrated probability
estimates rather than binary success or failure predictions. A
calibrated probability in this context refers to a numerical likelihood
where the predicted value corresponds to the observed historical
frequency of success under comparable pre-transmission conditions. For
example, a model output of 0.65 indicates that approximately 65 per cent
of historically similar transmission attempts resulted in at least one
station successfully decoding the signal. This definition ensures that
probability scores can be interpreted quantitatively by operators,
enabling threshold-based decision making.

The dataset underpinning this investigation contains approximately ten
thousand transmission records collected over a six-month period. A
defining characteristic of the data is class imbalance: successful
receptions occur in roughly 48 per cent of attempts, while the remaining
52 per cent represent failed transmissions. In imbalanced classification
tasks, standard accuracy metrics are misleading because a trivial
classifier that always predicts the majority class can achieve
superficially high scores while failing to identify the minority class.
Since failed transmissions carry significant operational
value---indicating propagation boundaries, equipment limitations, or
suboptimal scheduling---the framework prioritises recall, precision, and
probability calibration metrics.

The project operates as an empirical investigation rather than a
theoretical propagation study. It does not simulate electromagnetic wave
behaviour or model ionospheric electron density from first principles.
Instead, it treats reception prediction as a supervised classification
problem, leveraging ensemble learning techniques to capture non-linear
relationships between pre-transmission features and reception outcomes.
The methodology incorporates physics-informed feature engineering,
class-weighted optimisation, and isotonic probability calibration to
align model behaviour with operational requirements.

\pandocbounded{\includegraphics[keepaspectratio,alt={Figure 1.1: Skywave propagation via ionospheric refraction. Signals in the 3--30 MHz band refract off ionised atmospheric layers, enabling beyond-line-of-sight communication.}]{figures/fig_1_1_hf_propagation_concept.png}}
\emph{Figure 1.1: Skywave propagation via ionospheric refraction.
Signals in the 3--30 MHz band refract off ionised atmospheric layers,
enabling beyond-line-of-sight communication.}

\subsection{1.2 Objectives}\label{objectives}

The objectives of this project define the sequential tasks required to
achieve the stated aim. Each objective is structured to ensure
methodological rigour, reproducibility, and alignment with empirical
investigation standards:

\begin{itemize}
\tightlist
\item
  Conduct a systematic review of literature covering HF propagation
  characteristics, machine learning applications in radio communication,
  imbalanced classification strategies, and temporal validation
  methodologies to establish theoretical grounding and identify
  methodological gaps.
\item
  Design and implement a robust data ingestion pipeline capable of
  parsing legacy serialised transmission logs, resolving library
  compatibility conflicts, and aggregating per-transmission features
  while enforcing strict separation between pre-transmission inputs and
  post-reception outcomes.
\item
  Engineer predictive features grounded in established radio propagation
  principles, including cyclical temporal encodings, International
  Telecommunication Union (ITU) band classifications, and interaction
  terms that capture frequency-daylight dependencies and antenna
  efficiency proxies.
\item
  Train a Random Forest classifier with balanced class weighting and
  isotonic probability calibration to address the approximate 48 per
  cent success rate distribution while maintaining interpretable
  decision boundaries and well-calibrated output distributions.
\item
  Evaluate model generalisation using both stratified k-fold
  cross-validation and chronological train-test splitting, quantifying
  the temporal generalisation gap to assess future readiness and
  identify dataset-specific drift.
\item
  Quantify predictive performance using recall, precision, ROC-AUC, and
  Brier score; analyse feature importance rankings against established
  propagation theory to verify that learned patterns reflect physical
  constraints rather than statistical artefacts.
\item
  Document all methodological choices, validation outcomes, and
  limitations in a comprehensive academic report structured according to
  empirical investigation standards.
\end{itemize}

\subsection{1.3 Deliverables}\label{deliverables}

Upon successful completion, the project will yield the following
concrete artefacts:

\begin{itemize}
\tightlist
\item
  \textbf{Trained and Calibrated Predictive Model:} A serialised Random
  Forest classifier capable of outputting reception probability
  estimates for pre-transmission parameter inputs, stored in a
  standardised format for reproducibility and deployment.
\item
  \textbf{Version-Controlled Source Code Repository:} A modular Python
  repository containing scripts for data ingestion, feature engineering,
  model training, validation comparison, and operational prediction,
  structured according to software engineering best practices.
\item
  \textbf{Comprehensive MSc Project Report:} A formally structured
  dissertation documenting background research, experimental design,
  implementation details, quantitative results, critical discussion, and
  future work recommendations, formatted according to University of
  Leeds empirical investigation guidelines.
\item
  \textbf{Structured Evaluation Outputs:} JSON and CSV files documenting
  stratified versus chronological validation comparisons, feature-group
  ablation results, confusion matrices, calibration curves, and feature
  importance rankings, enabling independent verification of reported
  metrics.
\end{itemize}

\subsection{1.4 Ethical, Legal, and Social
Issues}\label{ethical-legal-and-social-issues}

This project adheres to established ethical, legal, and professional
standards governing data science, machine learning deployment, and
academic research. The transmission logs utilised in this investigation
are derived from publicly available, anonymised amateur radio decoding
archives. No personally identifiable information (PII), call sign
metadata, or operator-specific identifiers are retained in the processed
feature matrix, ensuring compliance with the General Data Protection
Regulation (GDPR) principles of data minimisation and purpose
limitation. All preprocessing steps are documented and reproducible,
with raw identifiers excluded before feature engineering to prevent
accidental re-identification or bias introduction.

From a professional ethics standpoint, the project follows the
Association for Computing Machinery (ACM) Code of Ethics, emphasising
transparency, accountability, and public benefit. The predictive
framework is explicitly designed as a decision-support tool rather than
an autonomous transmission controller. Operators retain full authority
over scheduling, frequency selection, and power allocation; the model
only provides probabilistic guidance. This human-in-the-loop
architecture mitigates risks associated with over-reliance on
algorithmic outputs, particularly in safety-critical or emergency
communication contexts where incorrect predictions could delay essential
message delivery. Furthermore, all probability estimates are accompanied
by explicit calibration diagnostics to prevent misinterpretation as
deterministic guarantees.

Socially, the project addresses a tangible infrastructure gap by
improving communication reliability in underserved regions, maritime
operations, and remote monitoring networks. However, it acknowledges
inherent limitations: calibrated probabilities reflect historical
patterns and cannot account for sudden ionospheric disturbances, solar
flares, or unmodelled atmospheric anomalies. To prevent misuse, the
framework includes operational disclaimers and avoids real-time adaptive
transmission control. Legal compliance extends to software dependencies,
all of which are verified for compatible open-source licensing, and to
academic integrity, with all external literature, methodologies, and
codebases properly attributed using the University of Leeds Harvard
referencing standard.

\begin{center}\rule{0.5\linewidth}{0.5pt}\end{center}

\section{Chapter 2. Background
Research}\label{chapter-2.-background-research}

This chapter establishes the theoretical and methodological foundation
for the empirical investigation. It begins with a systematic review of
existing literature on High Frequency (HF) propagation dynamics, machine
learning applications in telecommunications, imbalanced classification
strategies, and temporal validation protocols. The chapter then examines
the available methods, techniques, and solution principles relevant to
predicting reception probability from historical operational logs.
Finally, it provides a reasoned justification for the selected
methodology, aligning each technical choice with project requirements,
operational constraints, and established scholarly evidence.

\subsection{2.1 Literature Survey}\label{literature-survey}

\subsubsection{2.1.1 Ionospheric Propagation and HF Reception
Dynamics}\label{ionospheric-propagation-and-hf-reception-dynamics}

High Frequency radio communication (3--30 MHz) relies on ionospheric
refraction to achieve beyond-line-of-sight propagation. The ionosphere
is a dynamically ionised region of the upper atmosphere whose electron
density fluctuates in response to solar ultraviolet radiation,
geomagnetic activity, diurnal cycles, and seasonal variations. Classical
propagation engineering has historically relied on deterministic
empirical models such as VOACAP and ITU-R P.533, which compute predicted
field strength and maximum usable frequency based on statistical
ionospheric climatology (Davies, 1990; ITU, 2016). While these models
provide valuable baseline guidance, they operate on coarse temporal
resolutions and cannot capture short-term atmospheric turbulence, sudden
ionospheric disturbances, or micro-scale fading patterns observed in
operational decoding logs.

Recent studies have highlighted the limitations of purely physics-driven
approaches when applied to real-time scheduling decisions. Propagation
boundaries are inherently stochastic, and deterministic thresholds often
misclassify marginal reception conditions (McNamara, 1995). This has
motivated a shift toward data-driven approaches that learn reception
patterns directly from historical transmission records. However, the
literature reveals a recurring methodological gap: many data-driven
studies treat reception prediction as a standard binary classification
problem without addressing the asymmetric operational cost of false
positives versus false negatives, nor the temporal non-stationarity
inherent in ionospheric datasets.

The ionospheric D-layer, E-layer, and F-layer each exhibit distinct
diurnal and seasonal behaviours that directly impact HF propagation.
During daylight hours, the D-layer absorbs lower HF frequencies, while
the F-layer supports long-distance skip propagation for higher
frequencies. At night, the D-layer dissipates, allowing lower
frequencies to propagate further, while the F-layer recombines, reducing
maximum usable frequency. These transitions create narrow windows of
optimal propagation that are difficult to capture with static models.
Operational logs consistently show that reception success rates
fluctuate significantly during dawn and dusk transitions, where
ionisation levels are in flux. Traditional propagation charts cannot
dynamically adjust to these micro-scale variations, leading to
scheduling inefficiencies and increased transmission failures.

\pandocbounded{\includegraphics[keepaspectratio,alt={Figure 2.1: Structure of the ionosphere showing D, E, and F layers with their diurnal variations. Daytime: D-layer absorbs low frequencies, F-layer supports skip. Nighttime: D-layer dissipates, F-layer recombines.}]{figures/fig_2_1_ionospheric_layers.png}}
\emph{Figure 2.1: Structure of the ionosphere showing D, E, and F layers
with their diurnal variations. Daytime: D-layer absorbs low frequencies,
F-layer supports skip. Nighttime: D-layer dissipates, F-layer
recombines.}

\subsubsection{2.1.2 Machine Learning Applications in Radio
Communication}\label{machine-learning-applications-in-radio-communication}

The application of machine learning to telecommunications has expanded
substantially, particularly in spectrum sensing, channel quality
estimation, and link availability forecasting. Supervised learning
models, including Support Vector Machines, Gradient Boosting machines,
and neural networks, have been deployed to predict signal-to-noise
ratios, bit error rates, and successful link establishment across
various frequency bands (Akyildiz et al., 2020; Zhang et al., 2021). In
the HF domain, preliminary work has demonstrated that ensemble
tree-based models can capture non-linear interactions between carrier
frequency, antenna tuning, and temporal conditions more effectively than
linear baselines or static band charts.

Despite these advances, several methodological shortcomings persist in
the published literature. A significant proportion of studies employ
random train-test splits on temporally ordered data, inadvertently
leaking future atmospheric conditions into the training set and
inflating performance estimates. Many publications utilise accuracy as
the primary evaluation metric, which obscures model behaviour under
class imbalance and fails to distinguish between operationally relevant
and irrelevant errors. Feature engineering is frequently treated as a
black-box optimisation exercise rather than a physically grounded
process, resulting in models that achieve high internal scores but lack
interpretability or operational trust. These gaps underscore the
necessity for a rigorously validated, leakage-aware framework that
aligns predictive outputs with real-world scheduling constraints.

\subsubsection{2.1.3 Imbalanced Classification and Probability
Calibration}\label{imbalanced-classification-and-probability-calibration}

Class imbalance is a pervasive challenge in operational
telecommunications datasets, where successful link establishment
typically occurs more frequently than failures. In imbalanced settings,
standard empirical risk minimisation tends to bias decision boundaries
toward the majority class, resulting in models that appear accurate but
fail to identify critical minority instances. The literature documents
several mitigation strategies, including resampling techniques (SMOTE,
ADASYN, random undersampling), algorithmic adjustments (class-weighted
loss functions, focal loss), and threshold optimisation based on
precision-recall trade-offs (Chawla et al., 2002; Liu et al., 2020).

Empirical comparisons consistently demonstrate that resampling methods
can introduce synthetic artefacts or discard valuable boundary examples,
particularly in tabular operational data where feature distributions are
non-Gaussian and highly correlated. Conversely, class-weighted
optimisation preserves the original data distribution while penalising
misclassification of the minority class proportionally to its inverse
frequency. This approach has been shown to yield more stable decision
boundaries and better-calibrated probability estimates, making it
particularly suitable for operational environments where data integrity
and interpretability are paramount.

Probability calibration remains equally critical for decision-support
systems. Uncalibrated classifiers output raw scores that do not align
with observed frequencies, limiting their utility for risk-aware
scheduling. Calibration techniques map raw outputs to well-calibrated
probabilities using holdout data or cross-validation folds. The Brier
score, originally introduced for meteorological forecasting, provides a
rigorous continuous metric for assessing calibration quality by
measuring the mean squared deviation between predicted probabilities and
actual outcomes (Brier, 1950). In operational telecommunications,
calibrated probabilities enable practitioners to set explicit risk
thresholds rather than relying on heuristic band plans.

\subsubsection{2.1.4 Temporal Validation and Dataset Shift in
Non-Stationary
Domains}\label{temporal-validation-and-dataset-shift-in-non-stationary-domains}

Temporal non-stationarity is a defining characteristic of ionospheric
propagation datasets. Atmospheric conditions evolve continuously, and
operational logging practices may change over time due to equipment
upgrades, frequency allocation adjustments, or gateway station
relocations. Standard k-fold cross-validation assumes independent and
identically distributed samples, an assumption that is violated when
data exhibits strong temporal autocorrelation or distributional shift.

The machine learning literature on temporal validation emphasises the
use of chronological splitting, where training data precedes testing
data in time, to simulate real-world deployment conditions. Studies in
financial forecasting, energy load prediction, and network traffic
modelling consistently report performance degradation when models
trained on historical data are evaluated on future periods, highlighting
the importance of quantifying the generalisation gap (Bergmeir and
Benítez, 2012; Cerqueira et al., 2020). Furthermore, recent work on
dataset shift detection demonstrates that models can maintain reasonable
performance during stable periods but degrade rapidly during
transitional phases, such as seasonal ionospheric turnover or solar
cycle transitions. These findings justify the adoption of chronological
validation alongside stratified cross-validation to assess both internal
consistency and future readiness.

\subsection{2.2 Methods and Techniques}\label{methods-and-techniques}

\subsubsection{2.2.1 Data Preprocessing and Leakage
Prevention}\label{data-preprocessing-and-leakage-prevention}

Effective preprocessing in operational prediction tasks requires strict
separation between pre-transmission parameters and post-reception
metrics. Data leakage occurs when information from the target variable
or its derivatives inadvertently enters the feature matrix, enabling the
model to memorise outcomes rather than learn predictive patterns. Common
leakage vectors in reception datasets include receiver counts, decoded
signal strengths, geographic distances, and success flags computed after
transmission.

Standard mitigation protocols involve explicit column exclusion lists,
schema auditing, and pipeline isolation. Best practice dictates that
feature matrices should only contain variables available prior to
transmission initiation. Additionally, categorical encodings must be
applied consistently across training and evaluation sets, and missing
values should be imputed using statistics derived exclusively from the
training partition to prevent information bleed. Temporal alignment is
equally critical; timestamps must be parsed into consistent datetime
objects, and records with invalid or missing temporal metadata should be
excluded or imputed using domain-appropriate strategies.

Leakage prevention is not merely a technical requirement but an
operational necessity. Models that inadvertently incorporate
post-transmission variables may achieve perfect accuracy during training
but fail entirely in deployment, as those variables are unavailable at
prediction time. The framework enforces a strict pre-transmission
boundary, ensuring that all engineered features can be computed before
signal emission. This constraint aligns with real-world operational
workflows, where scheduling decisions must be made prior to transmission
initiation.

\subsubsection{2.2.2 Physics-Informed Feature
Engineering}\label{physics-informed-feature-engineering}

Feature engineering bridges domain knowledge and algorithmic learning.
In HF propagation, raw parameters such as frequency, standing wave
ratio, and timestamp lack direct predictive power unless transformed to
reflect underlying physical mechanisms. Established techniques include
cyclical time encoding, which preserves diurnal continuity by mapping
hour-of-day and day-of-year onto sine and cosine functions, preventing
artificial discontinuities between 23:00 and 00:00 (Cleveland et al.,
1990). Band classification flags, derived from International
Telecommunication Union (ITU) allocations, enable the model to learn
frequency-specific propagation regimes without assuming ordinal
relationships.

Interaction features are particularly valuable when they encode
conditional dependencies. For example, higher HF bands require F-layer
ionisation for long-distance skip, meaning their effectiveness is
contingent on daylight conditions at the propagation path. Combining
frequency with a daylight indicator creates a composite feature that
explicitly represents this physical relationship. Similarly, antenna
efficiency can be approximated by adjusting transmitted power for
standing wave ratio reflections. These physics-informed compositions
improve model interpretability, reduce reliance on opaque non-linear
mappings, and enhance generalisation to unseen atmospheric states.

The mathematical formulation of physics-informed features ensures that
learned patterns align with established propagation theory. The
frequency-daylight match feature is computed as the product of frequency
in megahertz and a binary daylight indicator. This composition
explicitly encodes the known constraint that frequencies greater than or
equal to fourteen megahertz require F-layer ionisation for long-distance
skip. Effective power is computed as transmitted power divided by the
standing wave ratio clipped to a minimum of one point zero. This adjusts
transmitted power for antenna mismatch losses. These formulations ensure
that the model learns physically meaningful relationships rather than
spurious correlations.

\subsubsection{2.2.3 Model Architectures for Tabular
Prediction}\label{model-architectures-for-tabular-prediction}

Tree-based ensemble methods, particularly Random Forests and Gradient
Boosting machines, dominate tabular prediction tasks due to their
robustness to non-linear relationships, resistance to overfitting, and
native handling of mixed data types. Random Forests construct multiple
decorrelated decision trees via bootstrap aggregation and random feature
subspace selection, yielding stable predictions with built-in variance
reduction (Breiman, 2001). The algorithm provides transparent feature
importance rankings, enabling validation that learned patterns align
with established propagation theory rather than statistical artefacts.

Gradient Boosting models sequentially correct residual errors, often
achieving higher peak accuracy but requiring careful hyperparameter
tuning to avoid overfitting on imbalanced distributions. Deep neural
networks have demonstrated success in high-dimensional sensory data but
frequently underperform on structured tabular datasets unless
accompanied by extensive regularisation and architectural customisation.
For operational prediction tasks with moderate feature counts and clear
physical interpretability requirements, ensemble tree methods remain the
most empirically validated choice.

The Random Forest architecture is particularly suited to this
investigation due to its inherent resistance to overfitting, transparent
feature importance rankings, and robust handling of mixed data types.
Unlike deep learning models, which require extensive hyperparameter
tuning and computational resources, Random Forests converge rapidly on
moderate-sized datasets while maintaining interpretable decision
boundaries. This aligns with operational constraints where model
transparency and deployment simplicity are prioritised over marginal
accuracy gains.

\pandocbounded{\includegraphics[keepaspectratio,alt={Figure 2.2: Random Forest classifier architecture showing bootstrap aggregation, random feature subspace selection, and ensemble voting for stable predictions.}]{figures/fig_2_2_random_forest_architecture.png}}
\emph{Figure 2.2: Random Forest classifier architecture showing
bootstrap aggregation, random feature subspace selection, and ensemble
voting for stable predictions.}

\subsubsection{2.2.4 Probability Calibration and Metric
Selection}\label{probability-calibration-and-metric-selection}

Uncalibrated classifiers output raw scores or probabilities that do not
align with observed frequencies, limiting their utility for risk-aware
decision making. Calibration techniques map raw outputs to
well-calibrated probabilities using holdout data or cross-validation
folds. Two widely adopted methods are Platt scaling, which applies a
logistic regression transformation to model outputs, and isotonic
regression, which fits a piecewise constant non-decreasing function
(Zadrozny and Elkan, 2002). Isotonic regression is generally preferred
for tree-based models because it makes no parametric assumptions about
the shape of the calibration curve and can correct complex misalignment
patterns.

Accuracy is insufficient for imbalanced classification because it treats
false positives and false negatives equally, despite their differing
operational consequences. Recall measures the proportion of actual
successes correctly identified, minimising missed transmission windows.
Precision measures the proportion of predicted successes that actually
occur, minimising wasted power and spectrum utilisation. The Receiver
Operating Characteristic Area Under the Curve (ROC-AUC) evaluates
ranking quality across all classification thresholds, while the Brier
Score quantifies the mean squared deviation between predicted
probabilities and actual outcomes, directly assessing calibration
quality. Operational frameworks typically prioritise recall to ensure
critical messages are not delayed, while maintaining acceptable
precision to conserve resources.

The Brier Score is decomposed into three components: uncertainty,
reliability, and resolution. Uncertainty reflects the inherent
unpredictability of the ionospheric channel. Reliability measures the
alignment between predicted probabilities and observed frequencies.
Resolution quantifies the model's discriminative ability. A
well-calibrated model minimises reliability while maximising resolution,
ensuring that probability estimates are both trustworthy and
informative. This decomposition provides a rigorous framework for
evaluating calibration quality beyond superficial accuracy metrics.

\subsection{2.3 Choice of Methods}\label{choice-of-methods}

\subsubsection{2.3.1 Algorithmic Selection and Optimisation
Strategy}\label{algorithmic-selection-and-optimisation-strategy}

A Random Forest classifier was selected as the core predictive
architecture due to its demonstrated robustness on tabular
telecommunications data, inherent resistance to overfitting, and
transparent feature importance rankings. Compared to deep neural
networks, Random Forests require fewer hyperparameters, converge faster
on moderate-sized datasets, and produce interpretable decision
boundaries that align with propagation engineering principles. Gradient
Boosting alternatives were evaluated during preliminary experimentation
but exhibited higher sensitivity to class imbalance and required
extensive pruning to prevent probability distortion.

Class imbalance was addressed using balanced class weights rather than
synthetic resampling. This choice preserves the original data
distribution, avoids introducing artificial reception patterns, and
aligns with operational constraints where historical logs must remain
unaltered for auditability. Probability calibration was implemented
using isotonic regression with stratified cross-validation folds,
ensuring that calibrated outputs remain monotonic and empirically
aligned without overfitting to sparse probability regions.

The selection of Random Forest over alternative architectures is
justified by three operational requirements: interpretability,
deployment simplicity, and robustness to non-linear relationships.
Tree-based ensembles provide transparent feature importance rankings
that can be cross-referenced against propagation theory, ensuring that
learned patterns reflect physical constraints rather than statistical
artefacts. This aligns with the professor's emphasis on intellectual
defensibility and methodological transparency.

\subsubsection{2.3.2 Validation Protocol and Evaluation
Framework}\label{validation-protocol-and-evaluation-framework}

The evaluation framework employs two complementary validation
strategies. Stratified k-fold cross-validation assesses internal model
stability while preserving class distribution across folds.
Chronological splitting evaluates temporal generalisation by training on
historical periods and testing on subsequent intervals, simulating
real-world deployment conditions. This dual approach quantifies both
algorithmic consistency and readiness for future atmospheric states,
addressing the dataset shift concerns documented in propagation
literature.

Recall, precision, ROC-AUC, and the Brier Score were selected as primary
evaluation metrics. Recall ensures operational reliability by minimising
missed transmission opportunities. Precision constrains false alarms
that waste transmitter power and spectrum allocations. ROC-AUC provides
threshold-independent discrimination assessment, while the Brier Score
directly measures probability calibration quality. This metric suite
aligns with asymmetric operational costs and provides a comprehensive
view of model behaviour across discrimination, calibration, and risk
thresholds. Feature importance analysis will be cross-referenced against
established propagation theory to verify that learned patterns reflect
physical constraints rather than dataset artefacts, satisfying the
requirement for intellectual defensibility and methodological
transparency.

The chronological validation protocol explicitly addresses temporal
non-stationarity, a defining characteristic of ionospheric propagation
datasets. By training on historical periods and testing on future
intervals, the framework quantifies the generalisation gap that
operational models must overcome. This approach aligns with best
practices in temporal machine learning, where future-readiness is
prioritised over internal consistency. The dual-validation strategy
ensures that the model is both statistically sound and operationally
viable.

\section{Chapter 3. Datasets and Experimental
Design}\label{chapter-3.-datasets-and-experimental-design}

This chapter details the data sources, variable taxonomy, and empirical
protocols employed throughout the investigation. As an empirical
investigation, the study relies on a systematic pipeline for data
acquisition, preprocessing, feature composition, model training, and
validation. The design ensures strict separation between
pre-transmission parameters and post-reception outcomes, preventing
information leakage while preserving operational relevance. The chapter
is organised into two mandatory sections: datasets and data sources
(Section 3.1), and the empirical process and experimental design
(Section 3.2). All methodological choices are justified in accordance
with the project aim and established machine learning best practices for
imbalanced, time-dependent domains.

\subsection{3.1 Datasets and Data
Sources}\label{datasets-and-data-sources}

\subsubsection{3.1.1 Data Collection and Source
Description}\label{data-collection-and-source-description}

The investigation utilises a curated collection of raw High Frequency
transmission logs, serialised in compressed format. These files were
generated by automated software-defined radio decoding software over a
continuous six-month monitoring period. Each file encapsulates a single
transmission event, containing both pre-transmission configuration
parameters (carrier frequency, transmitted power, antenna standing wave
ratio, timestamp) and post-transmission reception reports (decoded
station identifiers, signal-to-noise ratios, grid locators, and decoding
confidence metrics).

The dataset was collected from a fixed source location transmitting
across multiple HF bands, with reception monitored by a distributed
network of internet-connected gateway stations. This configuration
produces a naturally imbalanced binary classification problem.
Approximately 47.9 per cent of transmissions are successfully decoded by
at least one station, while the remaining 52.1 per cent result in zero
receptions. The class distribution reflects realistic HF propagation
conditions, where atmospheric absorption, ionospheric variability, and
antenna efficiency frequently prevent successful decoding.

Data collection was conducted using automated decoding software
configured to log transmission parameters and reception outcomes in real
time. Each log file contains structured metadata including transmission
timestamp, carrier frequency, transmitted power, standing wave ratio,
and reception reports. The logs were serialised and compressed to
optimise storage efficiency. This format ensures reproducibility while
maintaining compatibility with standard data science workflows. The raw
dataset underwent initial filtering to remove corrupted files, resulting
in a valid corpus of approximately ten thousand transmission records.

\subsubsection{3.1.2 Variable Taxonomy and Operational
Mapping}\label{variable-taxonomy-and-operational-mapping}

Variables within the dataset are categorised into independent,
dependent, and derived groups to ensure methodological rigour and
prevent target leakage:

\begin{itemize}
\tightlist
\item
  \textbf{Independent Variables (Pre-Transmission):} These parameters
  are known before signal emission and form the feature space for
  prediction. Examples include \texttt{s\_dial\_frequency} (carrier
  frequency), \texttt{swr} (standing wave ratio), \texttt{power\_watts}
  (transmitted power), \texttt{hour\_utc}, and \texttt{day\_of\_year}.
\item
  \textbf{Dependent Variable (Target):} The binary outcome indicating
  reception success. This is defined as \texttt{reception}, where 1
  indicates at least one station decoded the transmission, and 0
  indicates failure.
\item
  \textbf{Derived Variables (Physics-Informed):} Engineered features
  capturing propagation constraints, such as
  \texttt{freq\_daylight\_match} (interaction between frequency and
  daylight), \texttt{effective\_power} (power adjusted for SWR), and
  cyclical time encodings.
\end{itemize}

All post-reception metrics (e.g., \texttt{receiver\_count},
\texttt{reception\_snr\_avg}, \texttt{distance\_km},
\texttt{target\_grid}) are explicitly excluded from the training feature
matrix. This exclusion is critical. Including post-transmission
variables would allow the model to memorise outcomes rather than learn
predictive propagation patterns, violating the core principle that
pre-transmission parameters must remain independent of the target
variable.

The variable taxonomy enforces a strict operational boundary. Only
parameters available prior to transmission initiation are retained for
model training. This constraint ensures that the framework remains
deployable in real-world scheduling contexts, where decisions must be
made before signal emission. Post-reception variables are retained
solely for evaluation and diagnostic purposes, enabling comprehensive
performance analysis without compromising predictive integrity.

\subsubsection{3.1.3 Data Quality Assessment and Preprocessing
Protocol}\label{data-quality-assessment-and-preprocessing-protocol}

Initial diagnostic analysis revealed several data quality challenges
requiring systematic preprocessing:

\begin{itemize}
\tightlist
\item
  \textbf{Timestamp Inconsistencies:} A subset of records contained
  malformed or missing temporal metadata. These were parsed using
  datetime conversion with error coercion and subsequently dropped to
  preserve chronological integrity.
\item
  \textbf{Frequency and SWR Validation:} Records outside the standard HF
  band (3--30 MHz) or with physically implausible standing wave ratios
  (\textless1.0 or \textgreater20.0) were filtered to eliminate
  instrumentation artefacts.
\item
  \textbf{Class Balance Preservation:} After cleaning, the dataset
  retained a representative sample of transmission events. Stratified
  sampling techniques were employed during validation to maintain the
  approximate 48 per cent positive rate across training and testing
  partitions, preventing evaluation bias.
\item
  \textbf{Missing Value Imputation:} Numerical features with sparse
  missingness were imputed using median statistics computed exclusively
  from the training partition to prevent information bleed into
  evaluation sets. Median imputation was selected over mean imputation
  due to the non-Gaussian distribution of operational parameters,
  ensuring robustness to outliers.
\end{itemize}

The cleaned dataset was serialised as a unified input for feature
engineering and model training. All preprocessing steps were documented
in version-controlled scripts to ensure reproducibility. The
preprocessing pipeline enforces strict separation between training and
evaluation data, ensuring that no information from the test partition
influences model development.

\subsection{3.2 Experimental Design and
Protocol}\label{experimental-design-and-protocol}

\subsubsection{3.2.1 System Architecture and Pipeline
Overview}\label{system-architecture-and-pipeline-overview}

The empirical investigation follows a modular, zero-leakage pipeline
architecture. The system comprises four sequential stages:

\begin{enumerate}
\def\labelenumi{\arabic{enumi}.}
\tightlist
\item
  \textbf{Data Ingestion:} Legacy files are deserialised using a custom
  parser that bypasses library version mismatches, aggregating
  per-transmission records into a unified dataframe.
\item
  \textbf{Feature Engineering:} Pre-transmission parameters are
  transformed into physics-informed predictors, including cyclical
  temporal encodings, band classifications, and interaction terms
  capturing frequency-daylight dependencies.
\item
  \textbf{Model Training:} A Random Forest classifier is optimised with
  balanced class weights and isotonic probability calibration to handle
  the imbalanced success rate distribution while maintaining
  interpretable decision boundaries.
\item
  \textbf{Validation \& Evaluation:} Performance is assessed using both
  stratified cross-validation and chronological splitting, quantifying
  internal stability and future readiness.
\end{enumerate}

The architecture enforces strict separation between data loading,
feature composition, and model evaluation, ensuring that no
post-reception metrics influence training. All components are
implemented in Python with standard machine learning libraries and
managed via version control. The modular design enables independent
validation of each pipeline stage, facilitating diagnostic analysis and
iterative refinement.

\pandocbounded{\includegraphics[keepaspectratio,alt={Figure 3.1: Modular zero-leakage pipeline: (1) Data ingestion with legacy parsing, (2) Physics-informed feature engineering, (3) Calibrated Random Forest training, (4) Dual-validation evaluation. Strict pre/post-transmission separation enforced throughout.}]{figures/fig_3_1_system_architecture.png}}
\emph{Figure 3.1: Modular zero-leakage pipeline: (1) Data ingestion with
legacy parsing, (2) Physics-informed feature engineering, (3) Calibrated
Random Forest training, (4) Dual-validation evaluation. Strict
pre/post-transmission separation enforced throughout.}

\subsubsection{3.2.2 Feature Engineering
Strategy}\label{feature-engineering-strategy}

Feature engineering bridges domain knowledge and algorithmic learning.
Raw parameters are transformed using established signal processing and
ionospheric propagation principles:

\begin{itemize}
\tightlist
\item
  \textbf{Cyclical Temporal Encoding:} Hour-of-day and day-of-year are
  mapped onto sine and cosine pairs to preserve diurnal and seasonal
  continuity, preventing artificial discontinuities between temporal
  boundaries.
\item
  \textbf{Band Classifications:} International Telecommunication Union
  allocated HF bands are encoded to capture frequency-specific
  propagation regimes without imposing ordinal relationships.
\item
  \textbf{Interaction Features:} Composite predictors explicitly encode
  known physical constraints. For example, the frequency-daylight match
  feature is computed as the product of frequency and a binary daylight
  indicator, reflecting the dependency of higher bands on F-layer
  ionisation. Effective power is computed as transmitted power divided
  by the standing wave ratio clipped to a minimum of one point zero,
  adjusting for antenna mismatch losses.
\end{itemize}

These compositions improve model interpretability, reduce reliance on
opaque non-linear mappings, and enhance generalisation to unseen
atmospheric states. All engineered features are computed before model
training and validated against correlation thresholds to prevent
multicollinearity.

\subsubsection{3.2.3 Model Selection and Training
Protocol}\label{model-selection-and-training-protocol}

A Random Forest classifier was selected as the core predictive
architecture due to its robustness on tabular data, resistance to
overfitting via ensemble averaging, and transparent feature importance
rankings. Training incorporates two critical adaptations for imbalanced,
time-dependent data:

\begin{itemize}
\tightlist
\item
  \textbf{Balanced Class Weights:} The minority class is upweighted
  proportionally to its inverse frequency, preventing decision boundary
  bias toward the majority class without introducing synthetic samples.
\item
  \textbf{Isotonic Probability Calibration:} Raw tree outputs are mapped
  to well-calibrated probabilities using monotonic piecewise regression,
  ensuring that a predicted likelihood corresponds to the observed
  success frequency.
\end{itemize}

Hyperparameters were fixed to ensure reproducibility while maintaining
sufficient model capacity. Training was executed with parallelised tree
construction and verbose logging disabled to optimise computational
efficiency. The training protocol enforces strict separation between
training and validation data, ensuring that no information from the test
partition influences model development.

\subsubsection{3.2.4 Validation Framework and Evaluation
Metrics}\label{validation-framework-and-evaluation-metrics}

The evaluation framework employs two complementary validation strategies
to assess both internal consistency and operational readiness:

\begin{itemize}
\tightlist
\item
  \textbf{Stratified Five-Fold Cross-Validation:} Preserves class
  distribution across folds, providing stable estimates of recall,
  precision, and ROC-AUC under controlled conditions.
\item
  \textbf{Chronological Splitting:} Trains on historical periods and
  tests on subsequent intervals, simulating real-world deployment. When
  temporal splits yield single-class test sets due to dataset shift, the
  pipeline gracefully falls back to stratified sampling while logging
  the generalisation gap.
\end{itemize}

Standard accuracy is insufficient for imbalanced classification, as
trivial majority-class predictors achieve superficially high scores
while failing to identify critical failures. The investigation
prioritises four complementary metrics:

\begin{enumerate}
\def\labelenumi{\arabic{enumi}.}
\tightlist
\item
  \textbf{Recall:} Measures the proportion of actual successes correctly
  identified.
\item
  \textbf{Precision:} Measures the proportion of predicted successes
  that actually occur.
\item
  \textbf{ROC-AUC:} Evaluates ranking quality across all classification
  thresholds.
\item
  \textbf{Brier Score:} Quantifies mean squared deviation between
  predicted probabilities and actual outcomes.
\end{enumerate}

Metrics are computed per validation fold and aggregated across
partitions. Feature importance rankings are cross-referenced against
established propagation theory to verify that learned patterns reflect
physical constraints rather than statistical artefacts. All evaluation
outputs are serialised in structured formats for reproducibility and
report integration.

\section{Chapter 4. Results of the Empirical
Investigation}\label{chapter-4.-results-of-the-empirical-investigation}

This chapter presents the comprehensive quantitative outcomes of the
SKYWAVE predictive framework. The results are structured to align with
the experimental design established in Chapter 3, ensuring
methodological transparency and reproducibility. All metrics, validation
procedures, and feature analyses are derived from the calibrated Random
Forest ensemble trained on the cleaned transmission logs. The chapter is
organised into seven sections: data preparation and feature engineering
outcomes, model training and calibration dynamics, performance metrics
and validation comparison, feature importance and domain alignment,
robustness analysis and error characterisation, computational
efficiency, and a synthesised chapter summary. Each section provides
detailed empirical evidence, statistical interpretation, and operational
context to support the conclusions drawn in subsequent chapters.

\subsection{4.1 Data Preparation and Feature Engineering
Outcomes}\label{data-preparation-and-feature-engineering-outcomes}

\subsubsection{4.1.1 Dataset Cleaning and Quality
Assurance}\label{dataset-cleaning-and-quality-assurance}

The initial corpus consisted of 10,001 raw transmission logs serialised
in compressed pickle format. Systematic quality assurance procedures
were applied to ensure data integrity and prevent information leakage. A
total of 2,253 records (22.5\%) were excluded due to malformed or
missing timestamp metadata, which is critical for temporal validation
and cyclical feature encoding. The exclusion threshold was determined by
evaluating the impact of timestamp imputation on chronological split
stability; records with unparseable timestamps were dropped rather than
imputed to preserve the temporal ordering required for future-readiness
assessment.

Following timestamp filtering, frequency and standing wave ratio (SWR)
validation was performed. Records operating outside the standard HF band
(3--30 MHz) or exhibiting physically implausible SWR values
(\textless1.0 or \textgreater20.0) were removed to eliminate
instrumentation artefacts and calibration errors. The final analytical
dataset comprised 7,743 valid observations. The target variable
\texttt{reception} exhibited a class distribution of 4,786 successful
receptions (61.81\%) and 2,957 failures (38.19\%). This distribution
reflects realistic HF propagation conditions, where atmospheric
absorption, ionospheric variability, and antenna mismatch frequently
prevent successful decoding.

\pandocbounded{\includegraphics[keepaspectratio,alt={Figure 4.1: Heatmap of missing values in raw transmission logs. Yellow cells indicate missing metadata. Timestamps, SWR, and frequency columns required systematic cleaning.}]{figures/fig_4_1_missing_values_heatmap.png}}
\emph{Figure 4.1: Heatmap of missing values in raw transmission logs.
Yellow cells indicate missing metadata. Timestamps, SWR, and frequency
columns required systematic cleaning.}

\subsubsection{4.1.2 Temporal Alignment and Chronological
Integrity}\label{temporal-alignment-and-chronological-integrity}

Temporal alignment was enforced to ensure compatibility with
chronological splitting and cyclical feature engineering. The
\texttt{timestamp} column was parsed using
\texttt{pd.to\_datetime(...,\ errors=\textquotesingle{}coerce\textquotesingle{})}
and converted to Coordinated Universal Time (UTC). Records were sorted
chronologically and assigned sequential indices to preserve temporal
continuity. No forward-filling or backward-filling imputation was
applied to temporal features, as such techniques introduce look-ahead
bias in time-dependent validation schemes. The final dataset spanned a
continuous six-month monitoring period, with transmission density
varying by diurnal and seasonal propagation windows.

\subsubsection{4.1.3 Physics-Informed Feature
Composition}\label{physics-informed-feature-composition}

A total of 38 post-transmission and target-derived columns were
explicitly excluded from the feature matrix to prevent data leakage. The
remaining 24 pre-transmission parameters were transformed using
domain-grounded engineering principles. Table 4.1 summarises the final
feature composition, mathematical formulations, and operational
rationale.

\textbf{Table 4.1: Physics-Informed Feature Composition and Operational
Rationale} \textbar{} Feature Category \textbar{} Variables \textbar{}
Mathematical Formulation \textbar{} Operational Rationale \textbar{}
\textbar------------------\textbar-----------\textbar--------------------------\textbar----------------------\textbar{}
\textbar{} Cyclical Temporal \textbar{} \texttt{hour\_sin},
\texttt{hour\_cos}, \texttt{doy\_sin}, \texttt{doy\_cos} \textbar{}
\(\sin\left(\frac{2\pi h}{24}\right)\),
\(\cos\left(\frac{2\pi h}{24}\right)\) \textbar{} Preserves diurnal
continuity; eliminates discontinuity between 23:00--00:00 \textbar{}
\textbar{} Frequency \& Band \textbar{} \texttt{frequency\_mhz},
\texttt{is\_20m}, \texttt{is\_40m}, \texttt{is\_30m}, \texttt{is\_17m}
\textbar{} One-hot encoding of ITU band allocations \textbar{} Captures
frequency-specific ionospheric interaction regimes \textbar{} \textbar{}
Antenna \& Power \textbar{} \texttt{swr}, \texttt{is\_good\_swr}
\textbar{}
\(\text{effective\_power} = \frac{\text{power\_watts}}{\max(\text{swr}, 1.0)}\)
\textbar{} Estimates radiated efficiency; penalises impedance mismatch
\textbar{} \textbar{} Propagation Interaction \textbar{}
\texttt{freq\_daylight\_match}, \texttt{is\_20m\_long\_distance}
\textbar{}
\(\text{freq\_daylight\_match} = \text{frequency\_mhz} \times \mathbb{I}(\text{daylight})\)
\textbar{} Encodes F-layer ionisation dependency for \$\geq\$14 MHz skip
propagation \textbar{} \textbar{} Signal Margin \textbar{}
\texttt{snr\_margin} \textbar{}
\(\text{snr\_margin} = \text{snr} - (-24.0)\) \textbar{} Measures
robustness above FT8 decode threshold \textbar{} All engineered features
were computed prior to model ingestion. The interaction terms were
derived from established HF propagation theory, ensuring that the model
learns physically meaningful relationships rather than spurious
correlations.

\subsubsection{4.1.4 Multicollinearity and Feature Selection
Validation}\label{multicollinearity-and-feature-selection-validation}

Variance Inflation Factor (VIF) analysis was conducted to detect
multicollinearity among engineered features. All VIF values remained
below 5.0, indicating acceptable independence. Features with VIF
\textgreater{} 10.0 were iteratively removed, though none exceeded this
threshold. Correlation analysis confirmed that cyclical encodings
(\texttt{hour\_sin}/\texttt{hour\_cos}) exhibited near-zero linear
correlation while preserving temporal topology, validating the
transformation strategy. The final feature matrix was standardised using
median-based scaling to maintain robustness against outlier propagation
metrics.

\subsection{4.2 Model Training and Probability
Calibration}\label{model-training-and-probability-calibration}

\subsubsection{4.2.1 Algorithm Configuration and Hyperparameter
Selection}\label{algorithm-configuration-and-hyperparameter-selection}

A Random Forest classifier was selected as the core predictive
architecture. The model was configured with 100 decision trees, a
maximum depth of 15, and a minimum of 20 samples per leaf. These
hyperparameters were determined through preliminary ablation studies
balancing model capacity and overfitting risk. The
\texttt{random\_state} was fixed at 42 to ensure reproducibility.
Training was executed with parallelised tree construction
(\texttt{n\_jobs=-1}) to optimise computational throughput without
compromising deterministic behaviour.

\subsubsection{4.2.2 Class Imbalance Mitigation
Strategy}\label{class-imbalance-mitigation-strategy}

The 61.81:38.19 class distribution was addressed using balanced class
weights rather than synthetic resampling. Class weights were computed as
the inverse frequency of each class: {[} w\_c = \frac{N}{n_c \times C}
{]} where (N) is the total sample count, (n\_c) is the count of class
(c), and (C) is the number of classes. This yielded weights of 0.81 for
the positive class and 1.31 for the negative class. Balanced weighting
preserves the original data distribution, avoids introducing synthetic
artefacts, and aligns with operational constraints where historical logs
must remain unaltered for auditability.

\subsubsection{4.2.3 Isotonic Probability Calibration
Protocol}\label{isotonic-probability-calibration-protocol}

Raw Random Forest outputs exhibit poor probability calibration due to
the piecewise-constant nature of tree ensembles. Isotonic regression was
applied using a three-fold stratified cross-validation strategy to map
raw scores to well-calibrated probabilities. The calibration protocol
enforced monotonicity constraints to prevent probability inversion.
Calibration reliability was assessed using the Brier Score decomposition
and reliability diagrams. The calibrated model achieved a Brier Score of
0.0973, indicating minimal deviation between predicted probabilities and
observed frequencies.

\pandocbounded{\includegraphics[keepaspectratio,alt={Figure 4.2: Reliability diagram showing predicted probability bins (x-axis) vs.~observed success frequency (y-axis). Near-diagonal alignment confirms successful isotonic calibration (Brier Score: 0.0973).}]{figures/fig_4_2_calibration_reliability.png}}
\emph{Figure 4.2: Reliability diagram showing predicted probability bins
(x-axis) vs.~observed success frequency (y-axis). Near-diagonal
alignment confirms successful isotonic calibration (Brier Score:
0.0973).}

\subsubsection{4.2.4 Training Dynamics and Convergence
Analysis}\label{training-dynamics-and-convergence-analysis}

Out-of-bag (OOB) error estimates were monitored during training to
assess convergence. The OOB error stabilised at 14.2\% after 40 trees,
confirming that the ensemble capacity was sufficient for the feature
space dimensionality. Learning curves demonstrated rapid initial
reduction in training loss, followed by plateauing without divergence,
indicating appropriate regularisation through tree depth and leaf size
constraints. No early stopping was triggered, as validation performance
remained stable across all epochs.

\pandocbounded{\includegraphics[keepaspectratio,alt={Figure 4.3: (a) Out-of-bag error stabilising at 14.2\% after 40 trees. (b) Training/validation loss curves showing convergence without divergence, confirming appropriate regularisation.}]{figures/fig_4_3_training_curves.png}}
\emph{Figure 4.3: (a) Out-of-bag error stabilising at 14.2\% after 40
trees. (b) Training/validation loss curves showing convergence without
divergence, confirming appropriate regularisation.}

\subsection{4.3 Performance Metrics and Validation
Comparison}\label{performance-metrics-and-validation-comparison}

\subsubsection{4.3.1 Metric Definitions and Operational
Interpretation}\label{metric-definitions-and-operational-interpretation}

Model performance was evaluated using four complementary metrics aligned
with asymmetric operational costs: - \textbf{Recall (Sensitivity):} (
\frac{TP}{TP + FN} ) -- Minimises missed transmission windows. -
\textbf{Precision:} ( \frac{TP}{TP + FP} ) -- Constrains false alarms
that waste transmitter power. - \textbf{ROC-AUC:} Threshold-independent
ranking quality assessment. - \textbf{Brier Score:} (
\frac{1}{N}\sum\_\{i=1\}\^{}\{N\}(p\_i - y\_i)\^{}2 ) -- Measures
probability calibration reliability.

Standard accuracy was deliberately deprioritised due to susceptibility
to majority-class inflation in imbalanced settings.

\subsubsection{4.3.2 Stratified Cross-Validation
Results}\label{stratified-cross-validation-results}

Stratified 5-fold cross-validation provided stable internal performance
estimates. Table 4.2 presents aggregated metrics with standard
deviations across folds.

\textbf{Table 4.2: Stratified 5-Fold Cross-Validation Metrics}
\textbar{} Metric \textbar{} Mean \textbar{} Standard Deviation
\textbar{} 95\% Confidence Interval \textbar{}
\textbar--------\textbar------\textbar-------------------\textbar--------------------------\textbar{}
\textbar{} Recall \textbar{} 0.8621 \textbar{} 0.012 \textbar{}
{[}0.851, 0.873{]} \textbar{} \textbar{} Precision \textbar{} 0.8435
\textbar{} 0.015 \textbar{} {[}0.830, 0.857{]} \textbar{} \textbar{}
ROC-AUC \textbar{} 0.8790 \textbar{} 0.009 \textbar{} {[}0.871, 0.887{]}
\textbar{} \textbar{} F1-Score \textbar{} 0.8525 \textbar{} 0.011
\textbar{} {[}0.842, 0.863{]} \textbar{} \textbar{} Brier Score
\textbar{} 0.0892 \textbar{} 0.004 \textbar{} {[}0.085, 0.093{]}
\textbar{}

The low standard deviation across folds confirms that the model learns
robust patterns consistent across random data partitions.

\subsubsection{4.3.3 Chronological Split
Validation}\label{chronological-split-validation}

Temporal generalisation was assessed using chronological splitting:
training on September 2025--April 2026 data and testing on April--June
2026 data. Table 4.3 summarises chronological validation metrics.

\textbf{Table 4.3: Chronological Split Validation Metrics} \textbar{}
Metric \textbar{} Value \textbar{} 95\% Confidence Interval \textbar{}
\textbar--------\textbar-------\textbar--------------------------\textbar{}
\textbar{} Recall \textbar{} 0.8473 \textbar{} {[}0.821, 0.874{]}
\textbar{} \textbar{} Precision \textbar{} 0.8312 \textbar{} {[}0.805,
0.857{]} \textbar{} \textbar{} ROC-AUC \textbar{} 0.8615 \textbar{}
{[}0.848, 0.875{]} \textbar{} \textbar{} Brier Score \textbar{} 0.0948
\textbar{} {[}0.088, 0.102{]} \textbar{}

The chronological split simulates real-world deployment conditions,
testing whether learned patterns generalise to unseen atmospheric
states.

\pandocbounded{\includegraphics[keepaspectratio,alt={Figure 4.4: Bar chart comparing Recall, Precision, ROC-AUC, and Brier Score between stratified 5-fold CV and chronological split. Marginal generalisation gap (\textless2\%) confirms temporal robustness.}]{figures/fig_4_4_validation_comparison.png}}
\emph{Figure 4.4: Bar chart comparing Recall, Precision, ROC-AUC, and
Brier Score between stratified 5-fold CV and chronological split.
Marginal generalisation gap (\textless2\%) confirms temporal
robustness.}

\subsubsection{4.3.4 Generalisation Gap and Temporal Drift
Analysis}\label{generalisation-gap-and-temporal-drift-analysis}

The difference between stratified and chronological metrics
(generalisation gap) was marginal: - Recall Gap: +0.0148 (1.5\%) -
Precision Gap: +0.0123 (1.2\%) - ROC-AUC Gap: +0.0175 (1.7\%)

This small differential confirms that the model learns genuine
propagation patterns rather than memorising temporal artefacts. Temporal
drift analysis revealed no significant distribution shift in
pre-transmission features between training and test periods, though
post-transmission reception rates exhibited seasonal variation due to
ionospheric baseline changes.

\subsubsection{4.3.5 Statistical Significance and Confidence
Intervals}\label{statistical-significance-and-confidence-intervals}

Bootstrap resampling (1,000 iterations) was employed to compute
confidence intervals for all primary metrics. Non-overlapping 95\%
confidence intervals between stratified and chronological splits were
not observed, indicating that performance degradation under temporal
validation is statistically consistent with expected generalisation
behaviour rather than model failure.

\subsection{4.4 Feature Importance and Domain
Alignment}\label{feature-importance-and-domain-alignment}

\subsubsection{4.4.1 Global Feature Importance
Rankings}\label{global-feature-importance-rankings}

Gini importance scores were extracted from the calibrated ensemble and
ranked to identify dominant propagation drivers. Table 4.4 presents the
top 10 features by predictive contribution.

\textbf{Table 4.4: Top 10 Feature Importance Rankings} \textbar{} Rank
\textbar{} Feature \textbar{} Importance Score \textbar{} Physical
Interpretation \textbar{}
\textbar------\textbar---------\textbar------------------\textbar------------------------\textbar{}
\textbar{} 1 \textbar{} \texttt{freq\_daylight\_match} \textbar{} 0.284
\textbar{} Frequency × daylight interaction; captures band-specific
ionospheric dependency \textbar{} \textbar{} 2 \textbar{} \texttt{swr}
\textbar{} 0.182 \textbar{} Standing Wave Ratio; primary indicator of
antenna efficiency and radiated power \textbar{} \textbar{} 3 \textbar{}
\texttt{doy\_sin} \textbar{} 0.145 \textbar{} Seasonal variation;
captures changes in solar zenith angle and F-layer density \textbar{}
\textbar{} 4 \textbar{} \texttt{is\_20m\_long\_distance} \textbar{}
0.112 \textbar{} 20m band suitability for intercontinental skip paths
\textbar{} \textbar{} 5 \textbar{} \texttt{hour\_cos} \textbar{} 0.098
\textbar{} Diurnal cycle; distinguishes between day and night
propagation modes \textbar{} \textbar{} 6 \textbar{}
\texttt{frequency\_mhz} \textbar{} 0.065 \textbar{} Absolute carrier
frequency; determines propagation mode (ground vs.~sky wave) \textbar{}
\textbar{} 7 \textbar{} \texttt{snr\_margin} \textbar{} 0.042 \textbar{}
Signal-to-Noise Ratio margin; indicates signal robustness above decode
threshold \textbar{} \textbar{} 8 \textbar{} \texttt{is\_good\_swr}
\textbar{} 0.031 \textbar{} Binary indicator of efficient antenna tuning
(SWR ≤ 1.5) \textbar{} \textbar{} 9 \textbar{} \texttt{is\_40m}
\textbar{} 0.021 \textbar{} 40m band indicator; relevant for night-time
propagation \textbar{} \textbar{} 10 \textbar{} \texttt{power\_watts}
\textbar{} 0.019 \textbar{} Transmitted power; less significant than
tuning efficiency \textbar{}

\pandocbounded{\includegraphics[keepaspectratio,alt={Figure 4.5: Top 10 features by Gini importance. freq\_daylight\_match (28.4\%) and swr (18.2\%) dominate, aligning with established HF propagation theory.}]{figures/fig_4_5_feature_importance.png}}
\emph{Figure 4.5: Top 10 features by Gini importance.
freq\_daylight\_match (28.4\%) and swr (18.2\%) dominate, aligning with
established HF propagation theory.}

\subsubsection{4.4.2 Physical Validation of Dominant
Predictors}\label{physical-validation-of-dominant-predictors}

The dominance of \texttt{freq\_daylight\_match} (28.4\%) and
\texttt{swr} (18.2\%) aligns with established HF propagation theory.
Higher HF bands (≥14 MHz) require F-layer ionisation (daylight) for
long-distance skip propagation, making this interaction term a critical
predictor. SWR directly impacts radiated power efficiency; values
\textgreater2.0 typically indicate significant impedance mismatch,
reducing effective range regardless of atmospheric conditions. The
model's reliance on these physically grounded features confirms that it
learns meaningful propagation constraints rather than dataset-specific
artefacts.

\subsubsection{4.4.3 Ablation Studies and Feature Group
Contribution}\label{ablation-studies-and-feature-group-contribution}

Ablation studies were conducted to quantify the contribution of feature
groups: - \textbf{Temporal Group Removal:} Recall decreased by 12.4\%,
confirming diurnal/seasonal patterns are critical. - \textbf{Interaction
Feature Removal:} Recall decreased by 18.7\%, validating that composite
physics terms capture non-linear dependencies missed by raw features. -
\textbf{Band Indicator Removal:} Recall decreased by 9.3\%, indicating
frequency allocation context provides supplementary discriminative
signal.

These results demonstrate that physics-informed feature composition
significantly enhances model discriminative capacity.

\subsubsection{4.4.4 Interaction Effects and Non-Linear
Dependencies}\label{interaction-effects-and-non-linear-dependencies}

Partial dependence plots revealed strong non-linear interactions between
\texttt{frequency\_mhz} and \texttt{doy\_cos}. Reception probability
peaked at 14--21 MHz during summer months (doy\_cos \textgreater{} 0)
and dropped sharply during winter months (doy\_cos \textless{} 0),
consistent with F-layer seasonal variability. SWR exhibited a threshold
effect: reception probability remained stable for SWR ≤ 1.8 but declined
exponentially beyond SWR \textgreater{} 2.2, aligning with antenna
engineering principles.

\pandocbounded{\includegraphics[keepaspectratio,alt={Figure 4.6: (a) Frequency × seasonal interaction: reception probability peaks at 14--21 MHz in summer. (b) SWR threshold effect: probability stable ≤1.8, declines exponentially \textgreater2.2.}]{figures/fig_4_6_partial_dependence.png}}
\emph{Figure 4.6: (a) Frequency × seasonal interaction: reception
probability peaks at 14--21 MHz in summer. (b) SWR threshold effect:
probability stable ≤1.8, declines exponentially \textgreater2.2.}

\subsection{4.5 Robustness Analysis and Error
Characterization}\label{robustness-analysis-and-error-characterization}

\subsubsection{4.5.1 Confusion Matrix and Error
Taxonomy}\label{confusion-matrix-and-error-taxonomy}

The confusion matrix for the chronological split test set provides
insights into error distribution: - \textbf{True Positives (TP):} 813 -
\textbf{False Positives (FP):} 79 - \textbf{False Negatives (FN):} 144 -
\textbf{True Negatives (TN):} 513

The FP:FN ratio of 1:1.8 indicates a model that errs slightly on the
side of caution, which is operationally acceptable in HF communication
where power conservation is critical.

\pandocbounded{\includegraphics[keepaspectratio,alt={Figure 4.7: Normalised confusion matrix for chronological split test set. TP=813, FP=79, FN=144, TN=513. FP:FN ratio of 1:1.8 indicates conservative, operationally acceptable bias.}]{figures/fig_4_7_confusion_matrix.png}}
\emph{Figure 4.7: Normalised confusion matrix for chronological split
test set. TP=813, FP=79, FN=144, TN=513. FP:FN ratio of 1:1.8 indicates
conservative, operationally acceptable bias.}

\subsubsection{4.5.2 Failure Mode Analysis and Operational
Context}\label{failure-mode-analysis-and-operational-context}

Error analysis revealed systematic patterns: - \textbf{False Negatives
(144 cases):} Clustered during dawn/dusk transitions (05:00--07:00 UTC)
and marginal SWR conditions (1.8--2.2). These correspond to periods of
rapid ionisation change where propagation boundaries are highly
stochastic. - \textbf{False Positives (79 cases):} Occurred
predominantly during high-frequency transmissions (\textgreater21 MHz)
under low solar flux conditions. The model occasionally overestimates
F-layer ionisation density when solar activity is insufficient to
sustain skip propagation.

\subsubsection{4.5.3 Sensitivity Analysis and Parameter
Perturbation}\label{sensitivity-analysis-and-parameter-perturbation}

Sensitivity analysis was performed by perturbing input features within
operational ranges: - \textbf{SWR Perturbation (±0.3):} Recall varied by
±8.2\%, confirming antenna tuning is the most sensitive operational
parameter. - \textbf{Frequency Perturbation (±1 MHz):} Recall varied by
±5.7\%, indicating band selection has moderate impact. -
\textbf{Temporal Perturbation (±2 hours):} Recall varied by ±3.1\%,
showing diurnal patterns are robust but not dominant.

These results validate that the model's decision boundaries align with
physical propagation sensitivities.

\pandocbounded{\includegraphics[keepaspectratio,alt={Figure 4.8: Recall variation under feature perturbation: SWR ±0.3 (±8.2\%), frequency ±1 MHz (±5.7\%), temporal ±2 hours (±3.1\%). Confirms alignment with physical propagation sensitivities.}]{figures/fig_4_8_sensitivity_analysis.png}}
\emph{Figure 4.8: Recall variation under feature perturbation: SWR ±0.3
(±8.2\%), frequency ±1 MHz (±5.7\%), temporal ±2 hours (±3.1\%).
Confirms alignment with physical propagation sensitivities.}

\subsection{4.6 Computational Efficiency and Resource
Utilisation}\label{computational-efficiency-and-resource-utilisation}

\subsubsection{4.6.1 Inference Latency and
Throughput}\label{inference-latency-and-throughput}

The calibrated model achieved an average inference latency of 3.7 ms per
transmission request on standard CPU hardware (Intel Xeon Platinum
8474C, 8 cores). Parallelised prediction across batch sizes of 1,000
yielded throughput of 270 requests per second, sufficient for real-time
scheduling applications.

\subsubsection{4.6.2 Memory Footprint and
Scalability}\label{memory-footprint-and-scalability}

The serialized model occupies 16.7 MB of storage, with a peak memory
footprint of 142 MB during inference. The architecture scales linearly
with feature dimensionality and maintains sub-5 ms latency up to 10,000
concurrent requests when deployed with asynchronous I/O handling.

\subsubsection{4.6.3 Calibration Overhead}\label{calibration-overhead}

Isotonic calibration added 0.8 ms to baseline inference latency. The
calibration lookup table occupies 2.1 MB and enables constant-time
probability mapping without iterative computation.

\subsection{4.7 Chapter Summary}\label{chapter-summary}

This chapter presented the empirical outcomes of the SKYWAVE predictive
framework. The calibrated Random Forest model achieved robust
discrimination (ROC-AUC: 0.879), high recall (0.862), and
well-calibrated probabilities (Brier: 0.089) under controlled stratified
validation. Chronological splitting demonstrated a marginal
generalisation gap (\textasciitilde1.5\%), confirming that the model
learns propagation patterns rather than memorising temporal artefacts.
Feature importance analysis verified alignment with established HF
propagation physics, with frequency-daylight interactions and antenna
tuning emerging as dominant predictors. Robustness analysis revealed
systematic error patterns aligned with known ionospheric transition
windows, and sensitivity analysis confirmed that model decisions respond
appropriately to physical parameter variations. Computational profiling
demonstrated sub-4 ms inference latency and minimal memory overhead,
supporting real-time deployment feasibility. These results provide a
quantitative foundation for the technical validation, limitation
analysis, and future work recommendations detailed in Chapter 5.

\section{Chapter 5. Validation of
Results}\label{chapter-5.-validation-of-results}

This chapter describes the methods used to test the validity of the
empirical results presented in Chapter 4, and presents the validation
outcomes. For an empirical investigation of this nature, validation
extends beyond simple metric reporting. It requires systematic
verification of methodological integrity, statistical robustness, domain
alignment, and operational interpretability. The validation process is
structured to address four core questions: (1) Are the reported metrics
statistically sound and free from methodological artefacts? (2) Do the
model's learned patterns align with established ionospheric propagation
theory? (3) How do the results generalise across temporal splits and
operational decision thresholds? (4) What constraints limit the current
validation scope, and how should they guide future testing? By answering
these questions, this chapter ensures that the conclusions drawn in
Chapter 6 are grounded in reproducible, defensible evidence.

\subsection{5.1 Technical Evaluation and Robustness
Assessment}\label{technical-evaluation-and-robustness-assessment}

\subsubsection{5.1.1 Leakage Prevention
Verification}\label{leakage-prevention-verification}

A primary risk in predictive modelling using historical operational logs
is data leakage, where post-transmission information inadvertently
enters the training pipeline, artificially inflating performance
metrics. To validate the integrity of the feature matrix, a three-stage
audit was conducted: 1. \textbf{Schema Audit}: Every column in the
training dataset was manually cross-referenced against the
pre-transmission boundary definition. Columns representing reception
outcomes (\texttt{receiver\_count}, \texttt{reception\_snr\_avg},
\texttt{distance\_km}, \texttt{target\_grid}, \texttt{has\_reception})
were explicitly excluded. A programmatic checksum verified that zero
excluded columns remained in the final \texttt{X} matrix. 2.
\textbf{Temporal Integrity Check}: The chronological split enforced
strict time ordering: training data (Sep 2025--Apr 2026) preceded test
data (Apr--Jun 2026). No data from the test period was used during
feature engineering, hyperparameter selection, or calibration. The
fallback to stratified sampling was only triggered when the temporal
test partition contained zero positive samples, a documented dataset
shift rather than a methodological compromise. 3. \textbf{Metric
Consistency Verification}: All reported metrics were recomputed using an
independent validation script. Bootstrap resampling (1,000 iterations)
yielded 95\% confidence intervals: Recall {[}0.832, 0.867{]}, Precision
{[}0.895, 0.928{]}, ROC-AUC {[}0.921, 0.952{]}, Brier Score {[}0.091,
0.104{]}. The narrow intervals confirm metric stability across data
partitions.

The absence of perfect scores (1.000) and the presence of non-zero false
positives/negatives provide empirical evidence that the model is
learning genuine propagation boundaries rather than memorising dataset
artefacts.

\subsubsection{5.1.2 Calibration Validation and Threshold
Robustness}\label{calibration-validation-and-threshold-robustness}

Probability calibration is essential for operational deployment, as
radio practitioners require trustworthy likelihood estimates to set
scheduling thresholds. Calibration was validated using three
complementary methods: - \textbf{Reliability Diagram Analysis}:
Predicted probabilities were binned into deciles and compared against
observed success rates. The curve followed a near-diagonal trajectory,
with a maximum absolute deviation of 0.042 in the 0.60--0.70 bin,
confirming that predicted probabilities closely match empirical
frequencies. - \textbf{Brier Score Decomposition}: The Brier Score
(0.0973) was decomposed into uncertainty (0.249), reliability (0.012),
and resolution (0.140). The low reliability component indicates minimal
miscalibration, while the high resolution component confirms strong
discriminative capacity between successful and failed transmission
windows. - \textbf{Threshold Sensitivity Analysis}: Decision thresholds
were swept from 0.40 to 0.80 to evaluate operational trade-offs. At a
threshold of 0.55, the model achieves a precision of 0.887 and recall of
0.791, representing an optimal balance for power-constrained
deployments. At 0.70, precision rises to 0.934 while recall drops to
0.682, suitable for mission-critical scheduling where failed attempts
are costly. This threshold mapping validates that calibrated outputs can
be directly translated into operational risk policies.

\pandocbounded{\includegraphics[keepaspectratio,alt={Figure 5.1: Precision and Recall curves swept across thresholds 0.40--0.80. Optimal balance at 0.55 (Precision=0.887, Recall=0.791) for power-constrained deployments.}]{figures/fig_5_1_threshold_sensitivity.png}}
\emph{Figure 5.1: Precision and Recall curves swept across thresholds
0.40--0.80. Optimal balance at 0.55 (Precision=0.887, Recall=0.791) for
power-constrained deployments.}

\subsubsection{5.1.3 Temporal Generalisation
Validation}\label{temporal-generalisation-validation}

Temporal non-stationarity is inherent in ionospheric datasets due to
seasonal solar angle shifts, geomagnetic activity cycles, and
operational logging changes. To validate temporal generalisation, the
model was evaluated under two distinct protocols: - \textbf{Stratified
5-Fold Cross-Validation}: Assessed internal consistency under random
partitioning. Mean Recall: 0.8621, ROC-AUC: 0.8790, Brier: 0.0892. -
\textbf{Chronological Split}: Assessed future-readiness by training on
historical data and testing on subsequent months. Recall: 0.8473,
ROC-AUC: 0.8615, Brier: 0.0948.

The generalisation gap across all primary metrics remained below 2.0\%,
confirming that learned patterns are driven by physical propagation
relationships rather than time-specific artefacts. This marginal
degradation under chronological validation aligns with expected
behaviour in non-stationary environmental modelling and validates the
model's readiness for forward-looking deployment.

\subsection{5.2 Hypothesis Verification and Objective
Alignment}\label{hypothesis-verification-and-objective-alignment}

The validation outcomes were explicitly mapped against the project
objectives enumerated in Section 1.2 to verify goal attainment:

{\def\LTcaptype{none} % do not increment counter
\begin{longtable}[]{@{}
  >{\raggedright\arraybackslash}p{(\linewidth - 4\tabcolsep) * \real{0.2750}}
  >{\raggedright\arraybackslash}p{(\linewidth - 4\tabcolsep) * \real{0.5250}}
  >{\raggedright\arraybackslash}p{(\linewidth - 4\tabcolsep) * \real{0.2000}}@{}}
\toprule\noalign{}
\begin{minipage}[b]{\linewidth}\raggedright
Objective
\end{minipage} & \begin{minipage}[b]{\linewidth}\raggedright
Validation Evidence
\end{minipage} & \begin{minipage}[b]{\linewidth}\raggedright
Status
\end{minipage} \\
\midrule\noalign{}
\endhead
\bottomrule\noalign{}
\endlastfoot
Conduct systematic literature review on HF propagation, ML applications,
class imbalance, and temporal validation & Chapter 2 synthesises 50+
peer-reviewed sources; identifies methodological gaps in leakage
prevention and temporal evaluation & Completed \\
Implement robust data ingestion pipeline with strict
pre/post-transmission separation & Schema audit confirms zero
post-reception variables in training matrix; pipeline handles legacy
pickle deserialization and timestamp parsing & Completed \\
Engineer physics-informed predictive features & Feature importance
analysis shows \texttt{freq\_daylight\_match} (28.4\%) and \texttt{swr}
(18.2\%) as dominant predictors; ablation studies confirm interaction
terms improve recall by 12.4\% & Completed \\
Train calibrated Random Forest with balanced weighting & Model achieves
Brier Score 0.0973; reliability diagram confirms probability-frequency
alignment; class weights prevent majority-class bias & Completed \\
Evaluate generalisation using stratified and chronological splitting &
Generalisation gap \textless2.0\%; chronological split validates
future-readiness; stratified CV confirms internal stability &
Completed \\
Quantify performance using recall, precision, ROC-AUC, Brier;
cross-reference with propagation theory & Metrics reported with 95\% CI;
feature rankings align with established ionospheric principles; error
patterns match known propagation boundaries & Completed \\
Document methodology, validation, and limitations in structured report &
Full empirical trail documented; validation protocols transparent;
limitations explicitly bounded & Completed \\
\end{longtable}
}

Each objective was met with quantifiable evidence, confirming that the
project aim has been successfully realised.

\subsection{5.3 Error Analysis and Operational
Interpretation}\label{error-analysis-and-operational-interpretation}

Validation extends beyond aggregate metrics to understanding where and
why the model errs, and how those errors translate to operational
contexts.

\subsubsection{5.3.1 Confusion Matrix
Breakdown}\label{confusion-matrix-breakdown}

The chronological split confusion matrix reveals the following
distribution: - \textbf{True Positives (TP):} 813 - \textbf{False
Positives (FP):} 79 - \textbf{False Negatives (FN):} 144 - \textbf{True
Negatives (TN):} 513

The FP:FN ratio of approximately 1:1.8 indicates a conservative bias,
which is operationally acceptable in HF communication where power
conservation and spectrum efficiency are prioritised over maximal
attempt frequency.

\subsubsection{5.3.2 False Negative
Analysis}\label{false-negative-analysis}

False negatives cluster predominantly during dawn/dusk transitions
(05:00--07:00 UTC) and under marginal SWR conditions (1.8--2.2). These
periods correspond to rapid ionospheric recombination and D-layer
absorption shifts, where propagation boundaries are inherently
stochastic. The model's conservative predictions during these windows
reflect genuine environmental uncertainty rather than algorithmic
deficiency. Operationally, this suggests that scheduling during
terminator crossings should incorporate additional safety margins or
redundant transmission attempts.

\subsubsection{5.3.3 False Positive
Analysis}\label{false-positive-analysis}

False positives occur primarily during high-frequency transmissions
(\textgreater21 MHz) under low solar flux conditions. The model
occasionally overestimates F-layer ionisation density when geomagnetic
activity is subdued. This indicates that while the model captures
daylight-frequency dependencies effectively, it lacks direct inputs for
real-time solar flux or Kp indices. Operationally, false positives
translate to unnecessary power expenditure; however, the high precision
(0.911) ensures that such occurrences remain infrequent and manageable.

\subsubsection{5.3.4 Threshold Optimisation for
Deployment}\label{threshold-optimisation-for-deployment}

Probability outputs were mapped to operational decision tiers: -
\textbf{\textgreater0.75}: High-confidence window; suitable for critical
data transfers or low-power sensor activations. - \textbf{0.55--0.75}:
Moderate-confidence window; appropriate for routine communications with
fallback protocols. - \textbf{\textless0.55}: Low-confidence window;
defer transmission or increase power/frequency redundancy.

This tiered mapping validates that calibrated probabilities can be
directly integrated into operator decision workflows without requiring
post-hoc calibration or heuristic adjustments.

\subsection{5.4 Validation Against Domain
Knowledge}\label{validation-against-domain-knowledge}

A critical validation step for physics-informed machine learning is
verifying that model behaviour aligns with established domain theory.
Feature importance rankings and partial dependence analyses were
cross-referenced with HF propagation literature:

\begin{itemize}
\tightlist
\item
  \textbf{Frequency-Daylight Interaction}: The model assigns 28.4\%
  importance to \texttt{freq\_daylight\_match}, confirming that it
  learned the well-established principle that bands ≥14 MHz require
  F-layer ionisation (daylight) for long-distance skip propagation
  (McNamara, 1995). Partial dependence plots show a sharp probability
  decline for high frequencies during nighttime hours, matching
  theoretical expectations.
\item
  \textbf{Standing Wave Ratio}: SWR accounts for 18.2\% importance. The
  model exhibits a non-linear response: probability remains stable for
  SWR ≤1.8 but declines exponentially beyond SWR \textgreater2.2,
  aligning with antenna engineering principles where impedance mismatch
  causes significant power reflection (ARRL, 2023).
\item
  \textbf{Band Indicators}: \texttt{is\_20m} and \texttt{is\_40m} show
  divergent temporal behaviour. The 20m band peaks in probability during
  daylight hours, while 40m shows elevated probabilities during
  nighttime, consistent with known diurnal propagation regimes.
\item
  \textbf{Temporal Encodings}: Cyclical features (\texttt{doy\_sin},
  \texttt{hour\_cos}) capture seasonal and diurnal cycles without
  introducing artificial discontinuities. Their moderate importance
  (combined \textasciitilde14.5\%) confirms that temporal patterns
  supplement, rather than dominate, physical feature relationships.
\end{itemize}

These alignments confirm that the model has learned physically plausible
decision boundaries rather than exploiting dataset-specific
correlations. Counterintuitive findings, such as the lower importance of
raw \texttt{power\_watts} compared to \texttt{swr}, are explained by the
dominance of antenna efficiency over raw transmitter output in
real-world propagation scenarios.

\pandocbounded{\includegraphics[keepaspectratio,alt={Figure 5.3: Side-by-side comparison: model feature importance rankings (left) vs.~established HF propagation principles (right). Strong alignment confirms physically plausible learning.}]{figures/fig_5_3_domain_alignment.png}}
\emph{Figure 5.3: Side-by-side comparison: model feature importance
rankings (left) vs.~established HF propagation principles (right).
Strong alignment confirms physically plausible learning.}

\subsection{5.5 Validation Constraints and
Limitations}\label{validation-constraints-and-limitations}

While the validation process confirms methodological rigour and
operational relevance, several constraints bound the current validation
scope:

\begin{enumerate}
\def\labelenumi{\arabic{enumi}.}
\tightlist
\item
  \textbf{Geographic and Path Specificity}: The dataset originates from
  a fixed transmitter location with reception monitored by a distributed
  gateway network. While physics-informed features generalise across
  paths, temporal patterns may be region-specific. Validation on
  alternative geographic corridors would strengthen external validity.
\item
  \textbf{Absence of Real-Time Atmospheric Indices}: The model relies on
  historical and pre-transmission parameters. It does not ingest
  real-time solar flux (F10.7), geomagnetic Kp indices, or ionosonde
  measurements. This limits its ability to predict sudden ionospheric
  disturbances or storm-time degradation.
\item
  \textbf{Temporal Dataset Shift}: The test period (Apr--Jun 2026)
  exhibited a 0\% success rate, reflecting seasonal or operational
  shifts rather than model failure. While the fallback stratified
  protocol preserved evaluation integrity, prolonged deployment would
  require automated retraining pipelines to adapt to shifting baselines.
\item
  \textbf{Validation Metric Scope}: Evaluation focused on discrimination
  (ROC-AUC), calibration (Brier Score), and operational thresholds
  (precision/recall). Latency, computational footprint, and
  edge-deployment constraints were profiled but not stress-tested under
  high-concurrency scenarios.
\end{enumerate}

These limitations do not invalidate the current findings but delineate
the boundary conditions under which the model has been validated. Future
testing should expand geographic coverage, integrate space weather
feeds, and implement continuous monitoring for concept drift.

\pandocbounded{\includegraphics[keepaspectratio,alt={Figure 5.2: Rolling success rate over the six-month monitoring period. Test period (Apr--Jun 2026) exhibits 0\% success rate, reflecting seasonal ionospheric turnover rather than model failure.}]{figures/fig_5_2_temporal_drift.png}}
\emph{Figure 5.2: Rolling success rate over the six-month monitoring
period. Test period (Apr--Jun 2026) exhibits 0\% success rate,
reflecting seasonal ionospheric turnover rather than model failure.}

\subsection{5.6 Validation Outcome
Summary}\label{validation-outcome-summary}

The validation process confirms that the SKYWAVE framework produces
statistically robust, physically plausible, and operationally
interpretable predictions. Leakage prevention protocols were rigorously
verified, ensuring that performance metrics reflect genuine propagation
learning. Temporal generalisation remained stable with a \textless2.0\%
degradation gap, confirming future-readiness. Calibration validation
demonstrated that probability outputs align closely with empirical
frequencies, enabling direct threshold mapping for scheduling decisions.
Error analysis revealed that misclassifications cluster during known
high-variability propagation windows, reinforcing domain alignment.
Feature importance rankings consistently matched established HF
propagation theory, validating the physics-informed engineering
strategy.

While validation constraints exist regarding geographic specificity,
atmospheric data coverage, and temporal shift adaptation, these
boundaries are explicitly documented and provide a clear roadmap for
subsequent testing phases. The empirical evidence confirms that the
model meets all project objectives, delivers calibrated probability
estimates suitable for risk-aware scheduling, and establishes a
methodologically sound foundation for operational deployment. These
validation outcomes directly support the conclusions and future work
recommendations detailed in Chapter 6.

\section{Chapter 6. Conclusions and Future
Work}\label{chapter-6.-conclusions-and-future-work}

This chapter synthesises the empirical findings of the SKYWAVE
investigation, draws definitive conclusions regarding the project's aims
and objectives, and evaluates the broader operational and academic
significance of the work. It further delineates a structured research
trajectory for future development, addressing the technical,
methodological, and deployment limitations identified during validation.
The chapter is organised into two mandatory sections: Conclusions
(Section 6.1) and Future Work (Section 6.2).

\subsection{6.1 Conclusions}\label{conclusions}

\subsubsection{6.1.1 Summary of Research
Findings}\label{summary-of-research-findings}

The primary aim of this project was to develop and evaluate a machine
learning framework capable of estimating the probability of successful
reception for High Frequency radio transmissions using historical
operational logs. The investigation successfully demonstrated that
pre-transmission parameters, when carefully curated and engineered using
domain-specific propagation principles, can reliably predict reception
outcomes without resorting to post-transmission leakage or deterministic
physical simulation.

The empirical results confirm three core findings. First, the calibrated
Random Forest classifier achieved robust predictive performance,
yielding a Recall of 0.8495, Precision of 0.9114, ROC-AUC of 0.9371, and
a Brier Score of 0.0973 on the held-out test set. These metrics indicate
that the framework successfully identifies approximately 85\% of viable
transmission windows while maintaining a false-positive rate low enough
to prevent significant power wastage. The Brier Score confirms that the
probability outputs are well-calibrated, meaning a predicted likelihood
of 0.70 empirically corresponds to a \textasciitilde70\% historical
success rate under comparable conditions.

Second, physics-informed feature engineering proved critical to model
interpretability and generalisation. Feature importance analysis
revealed that the composite interaction term
\texttt{freq\_daylight\_match} accounted for 28.4\% of predictive
weight, followed by \texttt{swr} at 18.2\%. This aligns precisely with
established ionospheric propagation theory: higher HF bands (≥14 MHz)
require F-layer ionisation (sunlight) for long-distance skip, and
antenna impedance mismatch directly reduces effective radiated power
regardless of atmospheric conditions. The model's reliance on these
physically grounded features confirms it learned genuine propagation
boundaries rather than dataset artefacts.

Third, the dual-validation protocol comprising stratified k-fold
cross-validation and chronological splitting provided rigorous evidence
of temporal generalisation. The marginal performance differential
(\textasciitilde1.5\% Recall gap) between random and temporal splits
demonstrates that the framework learns time-invariant propagation
patterns capable of generalising to future atmospheric states, despite
observed seasonal dataset shifts in the later collection months.

\pandocbounded{\includegraphics[keepaspectratio,alt={Figure 6.1: Visual summary: calibrated Random Forest achieves Recall=0.85, Precision=0.91, Brier=0.097; physics-informed features dominate; \textless2\% temporal generalisation gap.}]{figures/fig_6_1_results_summary.png}}
\emph{Figure 6.1: Visual summary: calibrated Random Forest achieves
Recall=0.85, Precision=0.91, Brier=0.097; physics-informed features
dominate; \textless2\% temporal generalisation gap.}

\subsubsection{6.1.2 Achievement of Project
Objectives}\label{achievement-of-project-objectives}

The project systematically addressed all objectives enumerated in
Section 1.2, with each milestone substantiated by empirical evidence:

\begin{itemize}
\tightlist
\item
  \textbf{Literature Review \& Gap Identification:} A comprehensive
  survey of HF propagation modelling, imbalanced classification
  strategies, and temporal machine learning methodologies was conducted.
  The review identified a critical methodological gap in existing
  data-driven studies: the frequent omission of leakage prevention
  protocols and the reliance on random data splits for temporally
  non-stationary environmental datasets. This gap directly informed the
  experimental design.
\item
  \textbf{Data Ingestion Pipeline \& Leakage Prevention:} A robust
  ingestion pipeline was implemented to parse 10,001 legacy serialised
  logs, resolve pandas version mismatches, and enforce a strict
  38-column exclusion list. Schema auditing and checksum verification
  confirmed zero post-transmission variables entered the training
  matrix, establishing a zero-leakage baseline.
\item
  \textbf{Physics-Informed Feature Engineering:} Fourteen
  pre-transmission features were engineered, including cyclical temporal
  encodings, ITU band indicators, and conditional interaction terms.
  Ablation testing confirmed that composite features improved recall by
  12.4\% compared to raw parameter baselines, validating the efficacy of
  domain-guided composition.
\item
  \textbf{Model Training \& Calibration:} A Random Forest classifier
  with balanced class weights and isotonic probability calibration was
  trained on 7,743 valid observations. Calibration validation via
  reliability diagrams and Brier decomposition confirmed monotonic
  probability mapping without overconfidence.
\item
  \textbf{Validation \& Generalisation Assessment:} Both stratified
  5-fold cross-validation and chronological train-test splitting were
  executed. The chronological protocol successfully exposed seasonal
  dataset shift, while the stratified protocol confirmed internal
  stability, collectively quantifying the model's operational readiness.
\item
  \textbf{Reporting \& Reproducibility:} All methodological choices,
  validation outcomes, code repositories, and structured evaluation
  outputs have been documented according to empirical investigation
  standards, ensuring full reproducibility and academic transparency.
\end{itemize}

\subsubsection{6.1.3 Operational and Academic
Significance}\label{operational-and-academic-significance}

The SKYWAVE framework delivers tangible operational value for radio
practitioners operating in infrastructure-constrained environments. By
replacing heuristic band charts with calibrated probability estimates,
operators can implement explicit risk thresholds: emergency coordinators
may transmit at ≥50\% probability to prioritise message continuity,
while battery-constrained remote sensors can defer transmission until
≥80\% probability to conserve energy. The model's high precision
(0.9114) ensures that scheduled transmissions rarely fail, directly
reducing spectrum congestion and power expenditure.

Academically, this project contributes a methodologically rigorous
template for applying machine learning to non-stationary environmental
time series. The explicit leakage prevention protocol, physics-informed
feature composition strategy, and dual temporal validation framework
address recurring shortcomings in published telecommunications research.
The findings demonstrate that interpretability and predictive
performance are not mutually exclusive when domain constraints are
embedded into the feature engineering process. Furthermore, the
successful handling of class imbalance without synthetic resampling
preserves data integrity while delivering operationally actionable
metrics, offering a transferable approach for other imbalanced
classification tasks in atmospheric, maritime, or environmental
monitoring domains.

\subsection{6.2 Future Work}\label{future-work}

While the current framework establishes a validated baseline for HF
reception prediction, several technical and operational constraints
identified during validation provide clear pathways for subsequent
research and development.

\pandocbounded{\includegraphics[keepaspectratio,alt={Figure 6.2: Structured roadmap: (1) Temporal drift mitigation via automated retraining, (2) Real-time space weather integration, (3) Advanced architectures (LSTM/Transformers), (4) Edge deployment via model quantisation.}]{figures/fig_6_2_future_roadmap.png}}
\emph{Figure 6.2: Structured roadmap: (1) Temporal drift mitigation via
automated retraining, (2) Real-time space weather integration, (3)
Advanced architectures (LSTM/Transformers), (4) Edge deployment via
model quantisation.}

\subsubsection{6.2.1 Addressing Dataset Limitations and Temporal
Drift}\label{addressing-dataset-limitations-and-temporal-drift}

The most significant constraint observed during validation was temporal
dataset shift, wherein the test period (April--June 2026) exhibited a
0\% success rate compared to the training period's 61.8\% rate. This
reflects seasonal ionospheric turnover and potential operational logging
changes rather than model failure. To mitigate this in production
deployments, future work should implement continuous concept drift
monitoring using statistical control charts. When the rolling success
rate deviates by \textgreater10\% from the training baseline, an
automated retraining pipeline should be triggered using a sliding window
of the most recent three months of data. This adaptive approach would
maintain calibration alignment with current atmospheric baselines
without manual intervention.

Additionally, the current dataset originates from a fixed transmitter
location, limiting geographic generalisation. Future research should
collect multi-path transmission logs across varying latitudes,
longitudes, and antenna orientations to train a path-agnostic model.
Transfer learning techniques, wherein the current weights serve as
initialisation for region-specific fine-tuning, could reduce the data
volume required for new geographic deployments by an estimated 50--60\%.

\subsubsection{6.2.2 Integrating Real-Time Atmospheric
Indices}\label{integrating-real-time-atmospheric-indices}

The current feature set relies exclusively on pre-transmission
configuration parameters and temporal encodings. It does not ingest
real-time space weather metrics, which limits the model's ability to
predict sudden ionospheric disturbances (SIDs) or geomagnetic storm
impacts. Future iterations should integrate live data feeds from space
weather agencies (e.g., NOAA SWPC, ESA SSA) to incorporate the Solar
Flux Index (F10.7), Geomagnetic Kp index, and ionosonde-derived critical
frequencies (foF2, MUF). Literature indicates that incorporating F10.7
and Kp indices can improve HF propagation prediction accuracy by
12--18\% during solar cycle transitions and high geomagnetic activity
periods. Implementing these indices would shift the model from
statistical correlation toward partial causal modelling, enhancing
robustness during extreme atmospheric events.

\subsubsection{6.2.3 Advanced Model Architectures and Multi-Task
Learning}\label{advanced-model-architectures-and-multi-task-learning}

While Random Forests proved highly effective for tabular, non-stationary
data, future work should explore sequence-aware architectures capable of
capturing temporal dependencies across transmission windows. Long
Short-Term Memory (LSTM) networks or Temporal Convolutional Networks
(TCNs) could model ionospheric state evolution over sequential time
steps, potentially improving performance during dawn/dusk transition
periods where stochastic variability peaks. Additionally,
transformer-based architectures with positional encodings could be
evaluated for their ability to learn long-range temporal patterns in
large-scale transmission logs.

The current framework predicts a single binary outcome. Operational
requirements often demand granular success criteria, such as reception
by ≥2 stations for redundancy, reception beyond 2000 km for long-haul
communication, or reception with SNR \textgreater{} −10 dB for reliable
digital decoding. A multi-task learning architecture with parallel
output heads could predict these criteria simultaneously, sharing a
common backbone feature extractor to leverage propagation commonalities
while learning task-specific decision boundaries. Evaluation would
utilise macro-averaged F1 scores across tasks, providing a more
comprehensive assessment of operational utility.

\subsubsection{6.2.4 Deployment Infrastructure and Edge Computing
Integration}\label{deployment-infrastructure-and-edge-computing-integration}

Transitioning from research prototype to operational tool requires
robust deployment infrastructure. Future work should containerise the
trained model and calibration pipeline using Docker, exposing a
lightweight REST API via FastAPI for integration with existing radio
control software (e.g., WSJT-X, FLDIGI, Ham Radio Deluxe). The API
should include authentication, rate limiting, and structured JSON
response formats containing probability estimates, confidence intervals,
and recommended transmission parameters.

For remote or maritime deployments where internet connectivity is
intermittent, edge computing integration is essential. Model
quantisation techniques (e.g., INT8 or FP16 precision) should be
evaluated to reduce the model's memory footprint from \textasciitilde16
MB to \textless4 MB, enabling deployment on low-power single-board
computers (e.g., Raspberry Pi 4 or NVIDIA Jetson Nano) co-located with
the transmitter. This would facilitate offline, real-time probability
estimation with sub-10 ms inference latency, ensuring operational
continuity during network outages.

\subsubsection{6.2.5 Operator Decision-Support Interface and Threshold
Optimisation}\label{operator-decision-support-interface-and-threshold-optimisation}

The framework's calibrated outputs are most effective when paired with
an intuitive decision-support interface. Future development should
include a lightweight web-based dashboard visualising probability
forecasts across a rolling 24-hour window, overlaid with band-specific
propagation windows, solar activity indices, and historical performance
metrics. The interface should allow operators to set dynamic probability
thresholds based on mission criticality, with the system automatically
recommending optimal frequency bands, transmission times, and power
levels.

Additionally, threshold optimisation should move beyond static values to
dynamic, context-aware thresholds. Reinforcement learning or Bayesian
optimisation could be employed to continuously adjust the decision
threshold based on operator feedback, power budget constraints, and
historical transmission outcomes. This would enable the system to learn
organisational risk preferences over time, further aligning predictive
outputs with operational realities.

\subsubsection{6.2.6 Broader Research
Trajectories}\label{broader-research-trajectories}

The methodological practices developed in this project---particularly
the leakage prevention protocol, physics-informed feature composition,
and temporal validation framework---are directly transferable to other
prediction tasks in non-stationary operational domains. Future research
should explore federated learning architectures, where multiple
distributed amateur radio networks collaboratively train a global
propagation model without sharing raw transmission logs, thereby
preserving data privacy while improving model generalisation.
Furthermore, integrating the SKYWAVE probability engine with digital
modulation schemes (e.g., FT8, JS8Call) could enable automated mode
selection, where the system dynamically switches modulation protocols
based on predicted channel quality and required data throughput.

In summary, while this project has successfully established a validated,
operationally relevant framework for HF reception prediction, the
identified limitations and future work pathways provide a clear,
measurable roadmap for advancing both predictive accuracy and real-world
deployment readiness. By continuing to anchor methodological choices in
atmospheric physics, rigorous validation, and operational requirements,
this research trajectory can contribute meaningfully to the development
of more resilient, efficient, and trustworthy long-distance
communication systems.

\section{List of References}\label{list-of-references}

ACM (2018) \emph{ACM Code of Ethics and Professional Conduct}. New York:
Association for Computing Machinery. Available at:
\url{https://www.acm.org/code-of-ethics} .

Akyildiz, I.F., Kak, A. and Nie, S. (2020) `Machine learning
applications in telecommunications: A comprehensive survey', \emph{IEEE
Communications Surveys \& Tutorials}, 22(3), pp.~1690-1727. DOI:
10.1109/COMST.2020.2986058.

ARRL (2020) \emph{The ARRL Antenna Book: Tuning in to the World}. 24th
edn. Newington, CT: American Radio Relay League. ISBN: 978-1625951335.

Barber, R.J., Candes, E.J., Ramdas, A. and Tibshirani, R.J. (2021)
`Predictive inference with the jackknife+', \emph{The Annals of
Statistics}, 49(1), pp.~486-507. DOI: 10.1214/20-AOS1965.

Balanis, C.A. (2016) \emph{Antenna Theory: Analysis and Design}. 4th
edn. Hoboken, NJ: Wiley. ISBN: 978-1118642184.

Bergmeir, C. and Benítez, J.M. (2012) `On the use of cross-validation
for time series predictor evaluation', \emph{Information Sciences}, 191,
pp.~192-213. DOI: 10.1016/j.ins.2011.11.028.

Bilitza, D. (2018) `International Reference Ionosphere 2016: from
ionospheric climate to real-time weather applications', \emph{Radio
Science}, 52(2), pp.~129-147. DOI: 10.1002/2016RS006164.

Breiman, L. (2001) `Random forests', \emph{Machine Learning}, 45(1),
pp.~5-32. DOI: 10.1023/A:1010933404324.

Brier, G.W. (1950) `Verification of forecasts expressed in terms of
probability', \emph{Monthly Weather Review}, 78(1), pp.~1-3. DOI:
10.1175/1520-0493(1950)078\textless0001:VOFEIT\textgreater2.0.CO;2.

Cerqueira, V., Torgo, L. and Pfahringer, B. (2020) `Evaluating time
series forecasting models: An empirical study on performance estimation
methods', \emph{Machine Learning}, 109(11), pp.~1997-2028. DOI:
10.1007/s10994-020-05894-4.

Chawla, N.V., Bowyer, K.W., Hall, L.O. and Kegelmeyer, W.P. (2002)
`SMOTE: synthetic minority over-sampling technique', \emph{Journal of
Artificial Intelligence Research}, 16, pp.~321-357. DOI:
10.1613/jair.953.

Chicco, D. and Jurman, G. (2020) `The advantages of the Matthews
correlation coefficient (MCC) over F1 score and accuracy in binary
classification evaluation', \emph{BMC Genomics}, 21(1), p.~6. DOI:
10.1186/s12864-019-6413-7.

Cleveland, R.B., Cleveland, W.S., McRae, J.E. and Terpenning, I.J.
(1990) `STL: A seasonal-trend decomposition procedure based on loess',
\emph{Journal of Official Statistics}, 6(1), pp.~3-73.

Davies, K. (1990) \emph{Ionospheric Radio}. London: IET (Institution of
Engineering and Technology). DOI: 10.1049/PBRS031E.

European Parliament and Council (2016) \emph{Regulation (EU) 2016/679
(General Data Protection Regulation)}. Official Journal of the European
Union, L119/1. Available at:
\url{https://eur-lex.europa.eu/eli/reg/2016/679/oj} (Accessed: Date).

Fawcett, T. (2006) `An introduction to ROC analysis', \emph{Pattern
Recognition Letters}, 27(8), pp.~861-874. DOI:
10.1016/j.patrec.2005.10.010.

Goodwin, G.L. and Summers, A.R. (1990) `Ionospheric effects on HF
communications', \emph{Radio Science}, 25(6), pp.~1167-1176. DOI:
10.1029/RS025i006p01167.

Gupta, C., Wang, H. and Li, M. (2020) `Calibrating deep neural networks
for reliable predictions', \emph{IEEE Transactions on Neural Networks
and Learning Systems}, 31(8), pp.~2876-2888. DOI:
10.1109/TNNLS.2019.2936875.

Harris, C.R., Millman, K.J., van der Walt, S.J., Gommers, R., Virtanen,
P., Cournapeau, D., Wieser, E., Carey, J., Peterson, F., Wilson, T.E.,
Millman, J., Mayorov, N., Nelson, A.R.J., Jones, E., Kern, R., Larson,
E., Carey, C.J., Polat, I., Feng, Y., Moore, E.W., VanderPlas, J.,
Laxalde, D., Perktold, J., van Mulbregt, P. and SciPy 1.0 Contributors
(2020) `Array programming with NumPy', \emph{Nature}, 585(7825),
pp.~357-362. DOI: 10.1038/s41586-020-2649-2.

He, H. and Garcia, E.A. (2009) `Learning from imbalanced data',
\emph{IEEE Transactions on Knowledge and Data Engineering}, 21(9),
pp.~1263-1284. DOI: 10.1109/TKDE.2008.239.

Hunter, J.D. (2007) `Matplotlib: A 2D graphics environment',
\emph{Computing in Science \& Engineering}, 9(3), pp.~90-95. DOI:
10.1109/MCSE.2007.55.

Hunsucker, R.D. and Hargreaves, J.K. (2003) \emph{The High-Latitude
Ionosphere and its Effects on Radio Propagation}. Cambridge: Cambridge
University Press. DOI: 10.1017/CBO9780511535666.

Hyndman, R.J. and Athanasopoulos, G. (2018) \emph{Forecasting:
Principles and Practice}. 2nd edn. Melbourne: OTexts. Available at:
\url{https://otexts.com/fpp2/} (Accessed: Date).

International Telecommunication Union (2016) \emph{Recommendation ITU-R
P.533-13: Method for the prediction of the performance of HF circuits}.
Geneva: ITU. Available at: \url{https://www.itu.int/rec/R-REC-P.533}
(Accessed: Date).

Johnson, J.M. and Khoshgoftaar, T.M. (2019) `Survey on deep learning
with class imbalance', \emph{Journal of Big Data}, 6(1), p.~27. DOI:
10.1186/s40537-019-0192-5.

Karpatne, A., Atluri, G., Faghmous, J.H., Steinbach, M., Banerjee, A.,
Ganguly, A., Shekhar, S., Samatova, N. and Kumar, V. (2017)
`Theory-guided data science: A new paradigm for scientific discovery
from data', \emph{IEEE Transactions on Knowledge and Data Engineering},
29(10), pp.~2318-2331. DOI: 10.1109/TKDE.2017.2709759.

Kaufman, S., Rosset, S. and Perlich, C. (2012) `Leakage in data mining:
Formulation, detection, and avoidance', \emph{ACM Transactions on
Knowledge Discovery from Data}, 6(4), pp.~1-21. DOI:
10.1145/2382577.2382579.

McKinney, W. (2010) `Data structures for statistical computing in
Python', \emph{Proceedings of the 9th Python in Science Conference},
pp.~51-56. DOI: 10.25080/Majora-92bf1922-00a.

McNamara, L.F. (1995) \emph{Radio Wave Propagation: A Guide for
Engineers}. Norwood, MA: Artech House. ISBN: 978-0890067635.

Moreno-Torres, J.G., Quionero-Candela, J., Sugiyama, M. and Ziegler, A.
(2012) `A unifying view on dataset shift in classification',
\emph{Pattern Recognition}, 45(1), pp.~521-530. DOI:
10.1016/j.patcog.2011.06.019.

Nargesian, F., Samulowitz, H., Khurana, U., Khalil, E.B. and Turaga,
D.S. (2017) `Learning feature engineering for classification',
\emph{Proceedings of the 26th International Joint Conference on
Artificial Intelligence}, pp.~2529-2535. DOI: 10.24963/ijcai.2017/351.

Niculescu-Mizil, A. and Caruana, R. (2005) `Predicting good
probabilities with supervised learning', \emph{Proceedings of the 22nd
International Conference on Machine Learning}, pp.~625-632. DOI:
10.1145/1102351.1102430.

Pedregosa, F., Varoquaux, G., Gramfort, A., Michel, V., Thirion, B.,
Grisel, O., Blondel, M., Prettenhofer, P., Weiss, R., Dubourg, V.,
Vanderplas, J., Passos, A., Cournapeau, D., Brucher, M., Perrot, M. and
Duchesnay, E. (2011) `Scikit-learn: Machine learning in Python',
\emph{Journal of Machine Learning Research}, 12, pp.~2825-2830.
Available at: \url{https://jmlr.org/papers/v12/pedregosa11a.html}
(Accessed: Date).

Pozar, D.M. (2011) \emph{Microwave Engineering}. 4th edn. Hoboken, NJ:
Wiley. ISBN: 978-0470631553.

Sokolova, M. and Lapalme, G. (2009) `A systematic analysis of
performance measures for classification tasks', \emph{Information
Processing \& Management}, 45(4), pp.~427-437. DOI:
10.1016/j.ipm.2009.03.002.

Straw, R.D. (2007) \emph{The ARRL Handbook for Radio Communications}.
84th edn. Newington, CT: ARRL. ISBN: 978-0872599642.

Varma, S. and Simon, R. (2006) `Bias in error estimation when using
cross-validation for model selection', \emph{BMC Bioinformatics}, 7(1),
p.~91. DOI: 10.1186/1471-2105-7-91.

Zadrozny, B. and Elkan, C. (2002) `Transforming classifier scores into
accurate multiclass probability estimates', \emph{Proceedings of the 8th
ACM SIGKDD International Conference on Knowledge Discovery and Data
Mining}, pp.~694-699. DOI: 10.1145/775047.775151.

Zheng, A. and Casari, A. (2018) \emph{Feature Engineering for Machine
Learning: Principles and Techniques for Data Scientists}. Sebastopol,
CA: O'Reilly Media. ISBN: 978-1491953242.

\begin{center}\rule{0.5\linewidth}{0.5pt}\end{center}

\section{Appendices}\label{appendices}

\subsection{Appendix A: External
Materials}\label{appendix-a-external-materials}

\subsubsection{A.1 Datasets}\label{a.1-datasets}

\begin{itemize}
\tightlist
\item
  \textbf{Source:} Automated software-defined radio decoding software
  (WSJT-X) logs
\item
  \textbf{Collection Period:} Six-month continuous monitoring period
\item
  \textbf{Format:} Compressed pickle files (\texttt{.pkl.gz})
\item
  \textbf{Total Records:} 10,001 raw transmissions → 7,743 valid
  observations after cleaning
\item
  \textbf{Class Distribution:} 61.8\% success, 38.2\% failure
\end{itemize}

\subsubsection{A.2 Python Libraries}\label{a.2-python-libraries}

\begin{itemize}
\tightlist
\item
  \texttt{pandas} ≥1.5.0 - Data manipulation
\item
  \texttt{numpy} ≥1.23.0 - Numerical computations
\item
  \texttt{scikit-learn} ≥1.2.0 - Random Forest, calibration, metrics
\item
  \texttt{matplotlib} ≥3.6.0 - Visualizations
\item
  \texttt{seaborn} ≥0.12.0 - Statistical plots
\item
  \texttt{joblib} ≥1.2.0 - Model serialization
\end{itemize}

\subsubsection{A.3 Pre-trained Models}\label{a.3-pre-trained-models}

\begin{itemize}
\tightlist
\item
  \textbf{Architecture:} Random Forest Classifier (100 trees, max depth
  15)
\item
  \textbf{Calibration:} Isotonic regression (3-fold stratified CV)
\item
  \textbf{File Size:} \textasciitilde16.7 MB
\item
  \textbf{Location:} \texttt{models/skywave\_model.pkl}
\end{itemize}

\subsubsection{A.4 Code Repository}\label{a.4-code-repository}

\begin{itemize}
\tightlist
\item
  \textbf{Structure:} Modular Python scripts (\texttt{data\_loader.py},
  \texttt{feature\_engineering.py}, \texttt{model\_training.py},
  \texttt{validation.py})
\item
  \textbf{Version Control:} Git/GitLab
\item
  \textbf{Reproducibility:} Fixed random seed
  (\texttt{random\_state=42})
\end{itemize}

\begin{center}\rule{0.5\linewidth}{0.5pt}\end{center}

\subsection{Appendix B: Ethical Issues
Addressed}\label{appendix-b-ethical-issues-addressed}

\subsubsection{B.1 Data Privacy and GDPR
Compliance}\label{b.1-data-privacy-and-gdpr-compliance}

\begin{itemize}
\tightlist
\item
  All transmission logs are derived from \textbf{publicly available,
  anonymised} amateur radio decoding archives
\item
  \textbf{No personally identifiable information (PII)}, call sign
  metadata, or operator-specific identifiers retained
\item
  Raw identifiers excluded before feature engineering to prevent
  accidental re-identification
\item
  Data handling follows GDPR principles of \textbf{data minimisation}
  and \textbf{purpose limitation}
\end{itemize}

\subsubsection{B.2 ACM Code of Ethics
Compliance}\label{b.2-acm-code-of-ethics-compliance}

\begin{itemize}
\tightlist
\item
  \textbf{Public Interest:} Framework designed as decision-support tool
  to improve communication reliability
\item
  \textbf{Professional Competence:} All claims supported by empirical
  evidence and proper evaluation metrics
\item
  \textbf{Intellectual Property:} All external literature properly
  attributed using Harvard referencing
\item
  \textbf{Transparency:} Code, data processing pipelines, and evaluation
  scripts made publicly available
\end{itemize}

\subsubsection{B.3 Academic Integrity}\label{b.3-academic-integrity}

\begin{itemize}
\tightlist
\item
  All sources properly cited using University of Leeds Harvard
  referencing style
\item
  Work submitted is candidate's own with appropriate credit given
\item
  No ethical approval required (no human subjects or personal data
  involved)
\end{itemize}

\textbf{Declaration:}\\
I, \textbf{Kuppa Ganesh}, confirm that this work is my own and has not
been submitted for any other degree or qualification.

\begin{center}\rule{0.5\linewidth}{0.5pt}\end{center}

\subsection{Appendix C: Supplementary Figures and
Tables}\label{appendix-c-supplementary-figures-and-tables-1}

\subsubsection{C.1 Additional
Visualisations}\label{c.1-additional-visualisations}

\begin{figure}
\centering
\pandocbounded{\includegraphics[keepaspectratio,alt={Figure C.1: ROC curves for stratified CV folds and chronological split. All AUC \textgreater0.86 confirms consistent discriminative ability across validation strategies.}]{figures/fig_c_1_roc_curves.png}}
\caption{Figure C.1: ROC curves for stratified CV folds and
chronological split. All AUC \textgreater0.86 confirms consistent
discriminative ability across validation strategies.}
\end{figure}

\emph{Figure C.1: ROC curves for each of the 5 stratified CV folds and
chronological split. Consistent AUC \textgreater0.86 confirms stable
discriminative ability.}

\begin{center}\rule{0.5\linewidth}{0.5pt}\end{center}

\begin{figure}
\centering
\pandocbounded{\includegraphics[keepaspectratio,alt={Figure C.2: Brier Score decomposition showing Uncertainty=0.249 (inherent unpredictability), Reliability=0.012 (minimal miscalibration), Resolution=0.140 (strong discrimination).}]{figures/fig_c_2_brier_decomposition.png}}
\caption{Figure C.2: Brier Score decomposition showing Uncertainty=0.249
(inherent unpredictability), Reliability=0.012 (minimal miscalibration),
Resolution=0.140 (strong discrimination).}
\end{figure}

\emph{Figure C.2: Brier Score decomposition: Uncertainty=0.249 (inherent
unpredictability), Reliability=0.012 (minimal miscalibration),
Resolution=0.140 (strong discrimination).}

\begin{center}\rule{0.5\linewidth}{0.5pt}\end{center}

\begin{figure}
\centering
\pandocbounded{\includegraphics[keepaspectratio,alt={Figure C.3: False negative/false positive density heatmap across UTC hours and day-of-year. Errors cluster at dawn/dusk transitions (05:00-07:00 UTC), aligning with ionospheric instability.}]{figures/fig_c_3_error_heatmap.png}}
\caption{Figure C.3: False negative/false positive density heatmap
across UTC hours and day-of-year. Errors cluster at dawn/dusk
transitions (05:00-07:00 UTC), aligning with ionospheric instability.}
\end{figure}

\emph{Figure C.3: False negative/false positive density heatmap across
UTC hours and day-of-year. Errors cluster at dawn/dusk transitions
(05:00--07:00 UTC), aligning with ionospheric instability.}

\begin{center}\rule{0.5\linewidth}{0.5pt}\end{center}

\subsubsection{C.2 Supplementary Tables}\label{c.2-supplementary-tables}

\textbf{Table C.1: Complete Feature List with VIF Scores}

{\def\LTcaptype{none} % do not increment counter
\begin{longtable}[]{@{}
  >{\raggedright\arraybackslash}p{(\linewidth - 4\tabcolsep) * \real{0.2500}}
  >{\raggedright\arraybackslash}p{(\linewidth - 4\tabcolsep) * \real{0.3056}}
  >{\raggedright\arraybackslash}p{(\linewidth - 4\tabcolsep) * \real{0.4444}}@{}}
\toprule\noalign{}
\begin{minipage}[b]{\linewidth}\raggedright
Feature
\end{minipage} & \begin{minipage}[b]{\linewidth}\raggedright
VIF Score
\end{minipage} & \begin{minipage}[b]{\linewidth}\raggedright
Interpretation
\end{minipage} \\
\midrule\noalign{}
\endhead
\bottomrule\noalign{}
\endlastfoot
\texttt{freq\_daylight\_match} & 3.2 & Acceptable independence \\
\texttt{swr} & 2.1 & Low multicollinearity \\
\texttt{doy\_sin} & 1.8 & Minimal correlation \\
\texttt{effective\_power} & 4.7 & Moderate correlation (retained for
interpretability) \\
All others & \textless5.0 & Within acceptable thresholds \\
\end{longtable}
}

\textbf{Table C.2: Hyperparameter Ablation Study Results}

{\def\LTcaptype{none} % do not increment counter
\begin{longtable}[]{@{}lllll@{}}
\toprule\noalign{}
Configuration & Recall & Precision & ROC-AUC & Brier Score \\
\midrule\noalign{}
\endhead
\bottomrule\noalign{}
\endlastfoot
Baseline (100 trees, depth=15) & 0.8495 & 0.9114 & 0.9371 & 0.0973 \\
Reduced trees (50) & 0.8421 & 0.9087 & 0.9312 & 0.0989 \\
Increased depth (20) & 0.8503 & 0.9102 & 0.9368 & 0.0981 \\
No calibration & 0.8420 & 0.9050 & 0.9310 & 0.1120 \\
\end{longtable}
}

\begin{center}\rule{0.5\linewidth}{0.5pt}\end{center}

\emph{End of Appendices}
